\documentclass[11pt,openany]{book}
\usepackage{pdenotes}
\hypersetup{pdftitle={The Analysis of Nonlinear PDEs},pdfauthor={Amer Rawashdeh}}

% To print each hint and solution right under its exercise (e.g. for a handout), uncomment:
% \showsolutionsinline

\begin{document}
\frontmatter
\begin{titlepage}\centering\vspace*{1.6in}
{\Huge\bfseries The Analysis of Nonlinear PDEs\par}\vspace{14pt}
{\Large Foundations and case studies\par}\vspace{40pt}
{\large Amer Rawashdeh\par}\vspace{8pt}
{\normalsize Draft, \today\par}\vfill
{\small \copyright\ Amer Rawashdeh. Licensed under Creative Commons Attribution 4.0 (CC BY 4.0):\\
you may share and adapt this work, with credit. \url{https://creativecommons.org/licenses/by/4.0/}\par}
\end{titlepage}

\chapter*{Introduction}
These pages explain how existence, uniqueness and regularity are proved for nonlinear partial differential equations, through complete worked examples. Textbooks teach the tools one at a time; papers use them without comment. Here each \textbf{case study} walks through a real argument, and every tool it uses links to a \textbf{concept page} with a precise statement, the intuition behind it, and a reference for the proof.

You can start from a case study and follow the links back, or read the concept pages in order along the reading path below.

\begin{rmkbox}{How to read this}
Part~I follows the reading path layer by layer: each layer uses only the layers before it.
Part~II contains the case studies; each one can be read on its own, following the links back into Part~I. Dashed boxes mark proofs still to be written.
\end{rmkbox}

\tableofcontents

\mainmatter
\part{Foundations}

\chapter{The language: function spaces}
% Weak derivatives and Sobolev spaces

\section{Weak derivatives and Sobolev spaces}\label{p:sobolev-spaces}
\status{stub}

\subsection*{The idea}

PDE solutions are rarely smooth enough for classical derivatives, but they can be integrated against smooth test functions. A weak derivative is defined by the integration-by-parts formula, and Sobolev spaces measure how many weak derivatives a function has and in which $L^p$ space they lie. Almost every estimate in PDE analysis is a bound on a Sobolev norm.

\subsection*{Statement}

\begin{defnbox}{Definition}
\index{weak derivative}\index{Sobolev space}%
Let $\Omega \subset \R^n$ be open and $u, v \in L^1_{\mathrm{loc}}(\Omega)$. Then $v = \partial^\alpha u$ weakly if $\int_\Omega u\, \partial^\alpha \varphi \,dx = (-1)^{|\alpha|}\int_\Omega v\,\varphi\,dx$ for all $\varphi \in C_c^\infty(\Omega)$. For $k \in \N$ and $1 \le p \le \infty$,

\[
      W^{k,p}(\Omega) = \{ u \in L^p(\Omega) : \partial^\alpha u \in L^p(\Omega) \text{ for all } |\alpha| \le k \},
      \qquad \norm{u}_{W^{k,p}} = \sum_{|\alpha|\le k} \norm{\partial^\alpha u}_{L^p}.
    \]

$H^k = W^{k,2}$, and $H^1_0(\Omega)$ is the closure of $C_c^\infty(\Omega)$ in $H^1(\Omega)$: the $H^1$ functions that vanish on $\partial\Omega$ in the trace sense.
\end{defnbox}

\begin{defnbox}{The dual space $H^{-1}$}
\index{H-1@$H^{-1}$ (dual of $H^1_0$)}%
$H^{-1}(\Omega)$ is the dual of $H^1_0(\Omega)$: the space of bounded linear functionals $f$ on $H^1_0(\Omega)$. Write $\langle f, v\rangle$ for the value of $f$ at $v$; this is the pairing between $H^{-1}$ and $H^1_0$. The norm is

\[
      \norm{f}_{H^{-1}} = \sup\big\{ \langle f, v\rangle : v \in H^1_0(\Omega),\ \norm{v}_{H^1} \le 1 \big\}.
    \]

Each $f \in L^2(\Omega)$ defines an element of $H^{-1}$ by $\langle f, v\rangle = \int_\Omega f v\,dx$, so $H^1_0 \subset L^2 \subset H^{-1}$, each space continuously embedded in the next; for $f \in L^2$ the pairing is the $L^2$ inner product. Here $H^1_0$ is \emph{not} identified with its own dual: $L^2$ is the space in the middle.
\end{defnbox}

\begin{thmbox}{Theorem (what elements of $H^{-1}$ look like)}
$f \in H^{-1}(\Omega)$ if and only if there are $f^0, f^1, \dots, f^n \in L^2(\Omega)$ with $\langle f, v\rangle = \int_\Omega \big(f^0 v + \sum_{i=1}^n f^i \partial_i v\big)\,dx$ for all $v \in H^1_0(\Omega)$. In the sense of distributions, $f = f^0 - \sum_{i=1}^n \partial_i f^i$. Moreover $\norm{f}_{H^{-1}}$ is the infimum of $\big(\int_\Omega \sum_{i=0}^n |f^i|^2\,dx\big)^{1/2}$ over all such representations.
\end{thmbox}

\begin{defnbox}{Homogeneous spaces}
\index{Sobolev space!homogeneous}%
$D^{k,r}(\Omega) = \{ u \in L^1_{\mathrm{loc}} : \norm{\nabla^k u}_{L^r} < \infty \}$ with seminorm $\norm{\nabla^k u}_{L^r}$, and $D^k = D^{k,2}$. These record only the top derivative. They are convenient for quantities like the velocity, where the estimates naturally control $\nabla u$ and $\nabla^2 u$ but not $u$ itself.
\end{defnbox}

\begin{rmkbox}{Facts to know}
$W^{k,p}$ is a Banach space; $H^k$ is a Hilbert space; for $1 \le p < \infty$ smooth functions are dense in $W^{k,p}(\Omega)$ when $\partial\Omega$ is Lipschitz; and for $1 < p < \infty$ the space is reflexive.
\end{rmkbox}

\subsection*{Why it is plausible}

The definition is forced: it is the only one that agrees with the classical derivative for smooth $u$ and makes sense when $u$ is merely integrable. A function like $|x|$ on $(-1,1)$ has weak derivative $\operatorname{sign}(x)$, but $\operatorname{sign}(x)$ has no weak derivative in $L^1$: a jump is the first thing weak differentiation cannot absorb.

$H^{-1}$ is where one derivative of an $L^2$ function lives: $H^1_0$ functions have one derivative more than $L^2$, and elements of $H^{-1}$ have one derivative less. For example, if $u \in H^1_0(\Omega)$ then $\Delta u \in H^{-1}$, since $\langle \Delta u, v\rangle = -\int_\Omega \nabla u \cdot \nabla v\,dx$ is bounded by $\norm{\nabla u}_{L^2}\norm{v}_{H^1}$. So for the heat equation $u_t = \Delta u$ with $u \in L^2(0,T;H^1_0)$, the time derivative $u_t$ lies in $L^2(0,T;H^{-1})$: exactly the setting of the continuity theorem on \pref{bochner-spaces}{spaces of time-dependent functions}. In one dimension the point evaluation $v \mapsto v(0)$ belongs to $H^{-1}(-1,1)$, because $H^1$ functions are continuous there; in two or more dimensions it does not.

\subsection*{Exercises}

\begin{exercise}
  Let $u(x) = |x|$ on $(-1,1)$. Show directly from the definition that its weak derivative is $\operatorname{sign}(x)$. Then show that $\operatorname{sign}(x)$ has no weak derivative in $L^1_{\mathrm{loc}}(-1,1)$.
  \begin{hint}
    For the first part, split $\int_{-1}^1 |x|\varphi'(x)\,dx$ at $0$ and integrate by parts on each side; the boundary terms at $0$ cancel. For the second, compute $\int_{-1}^1 \operatorname{sign}(x)\varphi'(x)\,dx = -2\varphi(0)$. If $v$ were a weak derivative, testing with $\varphi$ supported away from $0$ forces $v = 0$ a.e., and then $-2\varphi(0) = 0$ for every $\varphi$, which is false.
  \end{hint}
  \begin{solution}
    Let $\varphi \in C_c^\infty(-1,1)$. Splitting at $0$ and integrating by parts on each side,

    \[ \begin{aligned}
             \int_{-1}^0 (-x)\varphi'\,dx &= \big[-x\varphi\big]_{-1}^0 + \int_{-1}^0 \varphi\,dx = \int_{-1}^0\varphi\,dx, \\
             \int_0^1 x\varphi'\,dx &= \big[x\varphi\big]_0^1 - \int_0^1\varphi\,dx = -\int_0^1\varphi\,dx,
             \end{aligned} \]

    since $\varphi(\pm1) = 0$. Adding, $\int_{-1}^1|x|\varphi'\,dx = -\int_{-1}^1\operatorname{sign}(x)\varphi\,dx$, which is the definition of $|x|' = \operatorname{sign}(x)$ weakly.

    Next, $\int_{-1}^1\operatorname{sign}(x)\varphi'\,dx = -\int_{-1}^0\varphi' + \int_0^1\varphi' = -\varphi(0) - \varphi(0) = -2\varphi(0)$. Suppose $v \in L^1_{\mathrm{loc}}$ were the weak derivative, so $\int\operatorname{sign}(x)\varphi' = -\int v\varphi$ for all $\varphi$. For $\varphi$ supported in $(0,1)$ the left side is $\int_0^1\varphi' = 0$, so $\int v\varphi = 0$ for all such $\varphi$ and $v = 0$ a.e. on $(0,1)$; likewise on $(-1,0)$. Then $-2\varphi(0) = 0$ for every test function, which fails for any $\varphi$ with $\varphi(0) = 1$.
  \end{solution}
\end{exercise}

\begin{exercise}
  Let $B$ be the unit ball in $\R^n$, $1 \le p < n$, and $u(x) = |x|^{-s}$ with $s > 0$. Show that $u \in W^{1,p}(B)$ if and only if $s < \frac{n-p}{p}$. Conclude that in $\R^3$ there are unbounded functions in $H^1(B)$.
  \begin{hint}
    Away from $0$, $|\nabla u| = s|x|^{-s-1}$, and in polar coordinates $\int_B |x|^{-(s+1)p}\,dx = c\int_0^1 r^{n-1-(s+1)p}\,dr$, finite exactly when $(s+1)p < n$. To see that the pointwise gradient is the weak gradient, integrate by parts on $B\setminus B_\varepsilon$: the boundary term is at most $C\varepsilon^{n-1-s}$, which tends to $0$ because $s < \frac{n-p}{p} \le n-1$. In $\R^3$ with $p = 2$, any $s \in (0,\frac12)$ works.
  \end{hint}
  \begin{solution}
    In polar coordinates, with $\omega_{n-1}$ the area of the unit sphere,

    \[ \int_B |x|^{-sp}\,dx = \omega_{n-1}\int_0^1 r^{n-1-sp}\,dr, \qquad \int_B |\nabla u|^p\,dx = s^p\omega_{n-1}\int_0^1 r^{n-1-(s+1)p}\,dr, \]

    using $\nabla u = -s|x|^{-s-2}x$ and $|\nabla u| = s|x|^{-s-1}$ away from $0$. The first is finite iff $sp < n$, the second iff $(s+1)p < n$, i.e. $s < \frac{n-p}{p}$; the second condition implies the first.

    It remains to check that the pointwise gradient is the weak gradient when $s < \frac{n-p}{p}$. For $\varphi \in C_c^\infty(B)$ and $\varepsilon > 0$, the divergence theorem on $B\setminus B_\varepsilon$ gives

    \[ \int_{B\setminus B_\varepsilon} u\,\partial_i\varphi\,dx = -\int_{B\setminus B_\varepsilon}\partial_iu\,\varphi\,dx + \int_{\partial B_\varepsilon} u\,\varphi\,\nu_i\,dS, \]

    with $\nu$ the normal on the small sphere. The last term is at most $\norm{\varphi}_{L^\infty}\varepsilon^{-s}\,\omega_{n-1}\varepsilon^{n-1}$, which tends to $0$ because $s < \frac{n-p}{p} = \frac np - 1 \le n-1$. Since $u, \partial_iu \in L^p(B) \subset L^1(B)$, letting $\varepsilon \to 0$ gives $\int_B u\,\partial_i\varphi = -\int_B \partial_iu\,\varphi$. Conversely, if $s \ge \frac{n-p}{p}$, a weak gradient would agree a.e. with the classical one on $B\setminus\{0\}$, which is not in $L^p$.

    For $n = 3$, $p = 2$: $\frac{n-p}{p} = \frac12$, so $|x|^{-s} \in H^1(B)$ for every $s \in (0,\frac12)$, and these functions are unbounded near $0$.
  \end{solution}
\end{exercise}

\begin{exercise}
  (a) Show that every $f \in L^2(\Omega)$ satisfies $\norm{f}_{H^{-1}} \le \norm{f}_{L^2}$. (b) On $\Omega = (-1,1)$, show that $\delta: v \mapsto v(0)$ belongs to $H^{-1}(\Omega)$. (c) Write $\delta = f^0 - (f^1)'$ with $f^0, f^1 \in L^2(-1,1)$.
  \begin{hint}
    (a) $|\int f v| \le \norm{f}_{L^2}\norm{v}_{L^2} \le \norm{f}_{L^2}\norm{v}_{H^1}$. (b) For $v \in H^1_0(-1,1)$, $v(0) = \int_{-1}^0 v'(x)\,dx$, so $|v(0)| \le \norm{v'}_{L^2} \le \norm{v}_{H^1}$ by Cauchy--Schwarz. (c) The formula in (b) says $\langle \delta, v\rangle = \int f^1 v'$ with $f^1 = \mathbf{1}_{(-1,0)}$ and $f^0 = 0$.
  \end{hint}
  \begin{solution}
    (a) For $v \in H^1_0$, $|\langle f,v\rangle| = |\int fv| \le \norm{f}_{L^2}\norm{v}_{L^2} \le \norm{f}_{L^2}\norm{v}_{H^1}$. Taking the supremum over $\norm{v}_{H^1} \le 1$ gives $\norm{f}_{H^{-1}} \le \norm{f}_{L^2}$.

    (b) For $v \in C_c^\infty(-1,1)$, $v(0) = v(0) - v(-1) = \int_{-1}^0 v'\,dx$, so by Cauchy--Schwarz $|v(0)| \le 1\cdot\norm{v'}_{L^2(-1,0)} \le \norm{v}_{H^1}$. The same bound, applied to differences, shows $\sup|v| \le \sqrt2\,\norm{v'}_{L^2}$, so $H^1_0$-limits of test functions converge uniformly; hence every $v \in H^1_0(-1,1)$ is continuous, $v(0)$ makes sense, and $|v(0)| \le \norm{v}_{H^1}$ still holds. So $\delta$ is a bounded linear functional: $\delta \in H^{-1}$ with $\norm{\delta}_{H^{-1}} \le 1$.

    (c) The formula in (b) reads $\langle\delta, v\rangle = \int_{-1}^1 f^1 v'\,dx$ with $f^1 = \mathbf{1}_{(-1,0)}$ and $f^0 = 0$. In the form of the theorem, $\delta = f^0 - (f^1)'$; indeed, for test functions, $\langle -(f^1)', \varphi\rangle = \int f^1\varphi' = \varphi(0) - \varphi(-1) = \varphi(0)$.
  \end{solution}
\end{exercise}

\subsection*{Where to find a proof}

\begin{itemize}
\item L. C. Evans, \emph{Partial Differential Equations}, 2nd ed., AMS, 2010, §5.2--5.5; §5.9.1 for $H^{-1}$ and the theorem describing its elements.
\item H. Brezis, \emph{Functional Analysis, Sobolev Spaces and Partial Differential Equations}, Springer, 2011, Chapters 8--9.
\end{itemize}

\neededfor

% Spaces of time-dependent functions

\section{Spaces of time-dependent functions}\label{p:bochner-spaces}

\status{stub}
\prereqs{sobolev-spaces}

\subsection*{The idea}

For an evolution equation we view a solution $u(x,t)$ as a curve $t \mapsto u(t)$ in a function space $X$. The notation $L^p(0,T;X)$ and $C([0,T];X)$ records how regular $u$ is in $x$ (through $X$) and in $t$ (through $L^p$ or $C$) separately. Definitions of solutions are written in this language.

\subsection*{Statement}

\begin{defnbox}{Definition}
Let $X$ be a Banach space. $L^p(0,T;X)$ is the space of strongly measurable $u:(0,T)\to X$ with $\norm{u}_{L^p(0,T;X)} = \big(\int_0^T \norm{u(t)}_X^p\,dt\big)^{1/p} < \infty$ (essential supremum for $p=\infty$). $C([0,T];X)$ is the space of continuous $u:[0,T]\to X$ with the norm $\max_t \norm{u(t)}_X$. $u_t = v$ weakly if $\int_0^T u\,\varphi'\,dt = -\int_0^T v\,\varphi\,dt$ in $X$ for all $\varphi\in C_c^\infty(0,T)$.
\end{defnbox}

\begin{thmbox}{Theorem (gaining continuity in time)}
If $u \in L^2(0,T;H^1_0(\Omega))$ and $u_t \in L^2(0,T;H^{-1}(\Omega))$, then $u \in C([0,T];L^2(\Omega))$ after modification on a set of measure zero, and $t\mapsto\norm{u(t)}_{L^2}^2$ is absolutely continuous with $\frac{d}{dt}\norm{u(t)}_{L^2}^2 = 2\langle u_t, u\rangle$ for a.e. $t$, where $\langle\cdot,\cdot\rangle$ is the pairing between $H^{-1}$ and $H^1_0$ (see \pref{sobolev-spaces}{Sobolev spaces}).
\end{thmbox}

\subsection*{Why it is plausible}

The theorem is what makes energy estimates legitimate: the formula $\frac{d}{dt}\frac12\norm{u}^2 = \langle u_t, u\rangle$ is exactly the step "multiply by $u$ and integrate", and it needs only that $u$ and $u_t$ live in dual spaces. Having $u$ in $H^1$ and $u_t$ in $H^{-1}$ puts $u$, on average, halfway between them in time, which is $L^2$.

\subsection*{Exercises}

\begin{exercise}
  Let $X$ be a Banach space, $v \in L^1(0,T;X)$, and $u(t) = \int_0^t v(s)\,ds$. Show that $u \in C([0,T];X)$ and that $u_t = v$ weakly.
  \begin{hint}
    Continuity: $\norm{u(t) - u(s)}_X \le \int_s^t \norm{v}_X$, which tends to $0$ as $t \to s$. Weak derivative: for $\varphi \in C_c^\infty(0,T)$ write $\int_0^T u\,\varphi'\,dt = \int_0^T \int_0^t v(s)\,ds\,\varphi'(t)\,dt$ and exchange the order of integration (Fubini for Bochner integrals); the inner integral becomes $\int_s^T \varphi' = -\varphi(s)$.
  \end{hint}
  \begin{solution}
    \emph{Continuity.} For $s < t$, $\norm{u(t) - u(s)}_X = \norm{\int_s^t v(\tau)\,d\tau}_X \le \int_s^t \norm{v(\tau)}_X\,d\tau$. Since $\norm{v}_X \in L^1(0,T)$, the right side tends to $0$ as $t - s \to 0$ (absolute continuity of the integral). So $u$ is (uniformly) continuous into $X$.

    \emph{Weak derivative.} For $\varphi \in C_c^\infty(0,T)$,

    \[ \begin{aligned}
             \int_0^T u(t)\varphi'(t)\,dt &= \int_0^T\!\!\int_0^t v(s)\,ds\,\varphi'(t)\,dt = \int_0^T v(s)\Big(\int_s^T\varphi'(t)\,dt\Big)ds \\
             &= \int_0^T v(s)\big(\varphi(T) - \varphi(s)\big)ds = -\int_0^T v\varphi\,ds.
             \end{aligned} \]

    Exchanging the order is allowed because $\int_0^T\!\int_0^T \norm{v(s)}_X|\varphi'(t)|\,ds\,dt < \infty$ (Fubini holds for Bochner integrals). This is the definition of $u_t = v$.
  \end{solution}
\end{exercise}

\begin{exercise}
  Let $w_k$ be the Dirichlet eigenfunctions of $-\Delta$ on $\Omega$, orthonormal in $L^2$, with eigenvalues $\lambda_k$, and let $u(t) = \sum_k a_k e^{-\lambda_k t} w_k$ be the solution of the heat equation with $u(0) = \sum_k a_k w_k \in L^2$. Show that $\int_0^T \norm{\nabla u(t)}_{L^2}^2\,dt \le \frac12\norm{u(0)}_{L^2}^2$. So $u \in L^2(0,T;H^1_0)$ even when $u(0) \notin H^1_0$.
  \begin{hint}
    The $w_k$ are orthogonal in $H^1_0$ too, with $\norm{\nabla w_k}_{L^2}^2 = \lambda_k$. So $\norm{\nabla u(t)}_{L^2}^2 = \sum_k \lambda_k a_k^2 e^{-2\lambda_k t}$, and integrating in $t$ gives $\sum_k a_k^2 \frac{1 - e^{-2\lambda_k T}}{2}$.
  \end{hint}
  \begin{solution}
    For eigenfunctions, $(\nabla w_j, \nabla w_k) = (-\Delta w_j, w_k) = \lambda_j\delta_{jk}$, so the $w_k$ are orthogonal in $H^1_0$ with $\norm{\nabla w_k}_{L^2}^2 = \lambda_k$. For $t > 0$ the series for $u(t)$ converges in $H^1_0$, and

    \[ \norm{\nabla u(t)}_{L^2}^2 = \sum_k \lambda_k a_k^2 e^{-2\lambda_k t}. \]

    All terms are nonnegative, so integration and summation can be exchanged (monotone convergence):

    \[ \int_0^T \norm{\nabla u}_{L^2}^2\,dt = \sum_k \lambda_k a_k^2\int_0^T e^{-2\lambda_kt}\,dt = \sum_k a_k^2\,\frac{1 - e^{-2\lambda_kT}}{2}
             \le \frac12\sum_k a_k^2 = \frac12\norm{u(0)}_{L^2}^2. \]

    This is finite for every $u(0) \in L^2$, while $u(0) \in H^1_0$ would need $\sum_k\lambda_k a_k^2 < \infty$. The heat equation gains a derivative on average in time.
  \end{solution}
\end{exercise}

\begin{exercise}
  Suppose $u \in L^2(0,T;H^1_0(\Omega))$, $u_t \in L^2(0,T;H^{-1}(\Omega))$, and $u_t = \Delta u$. Use the theorem above to show $\norm{u(t)}_{L^2}^2 + 2\int_0^t \norm{\nabla u}_{L^2}^2\,ds = \norm{u(0)}_{L^2}^2$. Deduce that this solution is determined by $u(0)$.
  \begin{hint}
    From \pref{sobolev-spaces}{Sobolev spaces}, $\langle \Delta u, u\rangle = -\norm{\nabla u}_{L^2}^2$. Integrate $\frac{d}{dt}\norm{u}_{L^2}^2 = 2\langle u_t, u\rangle$ from $0$ to $t$. For uniqueness, apply the identity to the difference of two solutions, which is again a solution, with zero initial value.
  \end{hint}
  \begin{solution}
    The equation $u_t = \Delta u$ in $L^2(0,T;H^{-1})$ means $\langle u_t, v\rangle = -\int\nabla u\cdot\nabla v$ for $v \in H^1_0$ and a.e. $t$. By the theorem, $t \mapsto \norm{u(t)}_{L^2}^2$ is absolutely continuous with

    \[ \frac{d}{dt}\norm{u(t)}_{L^2}^2 = 2\langle u_t, u\rangle = -2\norm{\nabla u(t)}_{L^2}^2 \quad\text{for a.e. } t. \]

    Integrating from $0$ to $t$ gives $\norm{u(t)}_{L^2}^2 + 2\int_0^t\norm{\nabla u}_{L^2}^2 = \norm{u(0)}_{L^2}^2$.

    If $u_1, u_2$ are two such solutions with the same initial value, $w = u_1 - u_2$ is again one (the equation is linear) with $w(0) = 0$. The identity gives $\norm{w(t)}_{L^2}^2 \le 0$, so $w \equiv 0$.
  \end{solution}
\end{exercise}

\subsection*{Where to find a proof}

\begin{itemize}
\item L. C. Evans, \emph{Partial Differential Equations}, 2nd ed., AMS, 2010, §5.9 (the theorem is §5.9.2, Theorem 3).
\end{itemize}

\neededfor

% Weak and weak* convergence

\section{Weak and weak* convergence}\label{p:weak-convergence}

\status{stub}
\prereqs{sobolev-spaces}

\subsection*{The idea}

Estimates give bounds, and bounds give convergent subsequences only in weak topologies. Almost every existence proof ends with "extract a weakly convergent subsequence and show the limit solves the equation". The hard part is always that weak limits do not commute with nonlinear functions.

\subsection*{Statement}

\begin{defnbox}{Definition}
$u_n \weakto u$ weakly in a Banach space $X$ if $f(u_n)\to f(u)$ for every $f \in X^*$. If $X = Y^*$, $u_n \weakstarto u$ weak* if $u_n(y) \to u(y)$ for every $y \in Y$.
\end{defnbox}

\begin{thmbox}{Theorem (compactness in weak topologies)}
(i) In a reflexive Banach space, every bounded sequence has a weakly convergent subsequence. (ii) If $Y$ is a separable Banach space, every bounded sequence in $Y^*$ has a weak* convergent subsequence. (iii) Norms are weakly lower semicontinuous: $u_n \weakto u$ implies $\norm{u} \le \liminf \norm{u_n}$.
\end{thmbox}

Typical use: $L^p(0,T;X)$ is reflexive for $1< p< \infty$ and $X$ reflexive; $L^\infty(0,T;X)$ is the dual of $L^1(0,T;X^*)$ when $X$ is reflexive and separable, which is where the weak* limits of energy bounds live.

\subsection*{Why it is plausible}

The standard example: $u_n(x) = \sin(nx)$ on $(0,2\pi)$ converges weakly to $0$ in $L^2$ (it oscillates faster and faster, so its integral against any fixed function vanishes), but $u_n^2 \weakto \frac12 \ne 0^2$. Weak convergence keeps averages and loses oscillations, so $\lim f(u_n) \ne f(\lim u_n)$ in general. Every nonlinear term in an equation must be handled with an extra argument: some strong convergence, or a structure that survives weak limits.

\subsection*{Exercises}

\begin{exercise}
  Show that $u_n(x) = \sin(nx) \weakto 0$ in $L^2(0,2\pi)$ and that $u_n^2 \weakto \frac12$, but that $u_n$ has no strongly convergent subsequence.
  \begin{hint}
    For $\varphi \in C_c^\infty(0,2\pi)$, integrate by parts: $|\int \sin(nx)\varphi| \le \frac1n\int|\varphi'|$. Extend to all $\varphi \in L^2$ by density, using $\norm{u_n}_{L^2} = \sqrt\pi$ for every $n$. For the squares use $\sin^2(nx) = \frac12 - \frac12\cos(2nx)$. A strong limit would have to be the weak limit $0$, but the norms stay $\sqrt\pi$.
  \end{hint}
  \begin{solution}
    For $\varphi \in C_c^\infty(0,2\pi)$, integration by parts gives $\int \sin(nx)\varphi\,dx = \frac1n\int\cos(nx)\varphi'\,dx$, so $|\int\sin(nx)\varphi| \le \frac1n\norm{\varphi'}_{L^1} \to 0$. For general $f \in L^2$ and $\varepsilon > 0$, choose $\varphi$ with $\norm{f - \varphi}_{L^2} < \varepsilon$. Since $\norm{u_n}_{L^2}^2 = \int_0^{2\pi}\sin^2(nx)\,dx = \pi$,

    \[ \Big|\int u_n f\Big| \le \Big|\int u_n\varphi\Big| + \norm{u_n}_{L^2}\norm{f - \varphi}_{L^2} \le \Big|\int u_n\varphi\Big| + \sqrt\pi\,\varepsilon, \]

    so $\limsup|\int u_nf| \le \sqrt\pi\,\varepsilon$ for every $\varepsilon$: $u_n \weakto 0$. The same argument gives $\cos(2nx) \weakto 0$, and $\sin^2(nx) = \frac12 - \frac12\cos(2nx) \weakto \frac12$.

    If a subsequence converged strongly in $L^2$ to some $g$, it would also converge weakly to $g$, so $g = 0$; but its norms are all $\sqrt\pi$, not tending to $0$.
  \end{solution}
\end{exercise}

\begin{exercise}
  In a Hilbert space, prove from the definition that $u_n \weakto u$ implies $\norm{u} \le \liminf\norm{u_n}$. Give an example where the inequality is strict.
  \begin{hint}
    $\norm{u}^2 = \lim (u_n, u) \le \liminf \norm{u_n}\,\norm{u}$. The sequence of Exercise 1 is an example.
  \end{hint}
  \begin{solution}
    The map $v \mapsto (v, u)$ is a bounded linear functional, so $(u_n, u) \to (u, u) = \norm{u}^2$. By Cauchy--Schwarz, $(u_n, u) \le \norm{u_n}\,\norm{u}$, hence $\norm{u}^2 \le \liminf_n\norm{u_n}\,\norm{u}$. If $u = 0$ there is nothing to prove; otherwise divide by $\norm{u}$. For $u_n = \sin(nx)$ in $L^2(0,2\pi)$: $\norm{0} = 0 < \sqrt\pi = \liminf\norm{u_n}$.
  \end{solution}
\end{exercise}

\begin{exercise}
  In a Hilbert space, show that $u_n \weakto u$ together with $\norm{u_n} \to \norm{u}$ implies $u_n \to u$ strongly. This is the pattern behind every "weak convergence plus convergence of the norm" argument.
  \begin{hint}
    Expand $\norm{u_n - u}^2 = \norm{u_n}^2 - 2(u_n, u) + \norm{u}^2$ and let $n \to \infty$.
  \end{hint}
  \begin{solution}
    $\norm{u_n - u}^2 = \norm{u_n}^2 - 2(u_n, u) + \norm{u}^2 \to \norm{u}^2 - 2\norm{u}^2 + \norm{u}^2 = 0$, using $\norm{u_n} \to \norm{u}$ and $(u_n, u) \to \norm{u}^2$ (weak convergence tested against $u$).
  \end{solution}
\end{exercise}

\subsection*{Where to find a proof}

\begin{itemize}
\item H. Brezis, \emph{Functional Analysis, Sobolev Spaces and Partial Differential Equations}, Springer, 2011, Chapter 3.
\item L. C. Evans, \emph{Partial Differential Equations}, 2nd ed., AMS, 2010, Appendix D.4.
\end{itemize}

\neededfor


\chapter{The inequality toolkit}
% Hölder, Young, and Lᵖ interpolation

\section{Hölder, Young, and L\textsuperscript{\textit{p}} interpolation}\label{p:elementary-inequalities}
\status{stub}

\subsection*{The idea}

Three inequalities do most of the bookkeeping in every estimate: Hölder splits an integral of a product into norms, Young splits a product of numbers into a sum (so that one piece can be absorbed into the left-hand side), and interpolation bounds an intermediate $L^p$ norm by two others.

\subsection*{Statement}

\begin{thmbox}{Hölder}
If $1 \le p, q, r \le \infty$ with $\frac1p + \frac1q = \frac1r$, then $\norm{fg}_{L^r} \le \norm{f}_{L^p}\norm{g}_{L^q}$.
\end{thmbox}

\begin{thmbox}{Young}
If $1 < p, q < \infty$ with $\frac1p + \frac1q = 1$ and $a, b \ge 0$, then $ab \le \frac{a^p}{p} + \frac{b^q}{q}$. With a parameter: for every $\varepsilon > 0$, $ab \le \varepsilon a^p + C(\varepsilon) b^q$, and in particular $ab \le \varepsilon a^2 + \frac{1}{4\varepsilon} b^2$ ("Cauchy's inequality with $\varepsilon$").
\end{thmbox}

\begin{thmbox}{Interpolation of L\textsuperscript{\textit{p}} norms}
If $1 \le p \le r \le q \le \infty$ and $\frac1r = \frac{\eta}{p} + \frac{1-\eta}{q}$ with $\eta\in[0,1]$, then $\norm{f}_{L^r} \le \norm{f}_{L^p}^{\eta}\norm{f}_{L^q}^{1-\eta}$.
\end{thmbox}

\subsection*{Why it is plausible}

The pattern to recognize: a bad term $\int |\nabla u|\,|F|$ is bounded by Hölder, then by $\varepsilon\norm{\nabla u}_{L^2}^2 + C(\varepsilon)\norm{F}_{L^2}^2$ by Young, and the first piece is moved to the left-hand side, where the dissipation $\mu\norm{\nabla u}_{L^2}^2$ absorbs it. Interpolation is how a conserved low norm (like the mass $\norm{\rho}_{L^1}$) and an energy norm (like $\norm{\rho}_{L^\gamma}$) together control the norm a term actually needs; the exponent $\eta$ is fixed by matching the reciprocals, which is just a weighted average.

\subsection*{Exercises}

\begin{exercise}
  Prove Young's inequality $ab \le \frac{a^p}{p} + \frac{b^q}{q}$ from the concavity of $\log$. Then deduce $ab \le \varepsilon a^2 + \frac{1}{4\varepsilon}b^2$.
  \begin{hint}
    For $a, b > 0$: $\log\big(\frac1p a^p + \frac1q b^q\big) \ge \frac1p\log a^p + \frac1q\log b^q = \log(ab)$. For the second inequality, apply the case $p = q = 2$ to the product $(\sqrt{2\varepsilon}\,a)(b/\sqrt{2\varepsilon})$.
  \end{hint}
  \begin{solution}
    If $a = 0$ or $b = 0$ there is nothing to prove. Otherwise, $\log$ is concave on $(0,\infty)$, so with $\theta = \frac1p$, $1 - \theta = \frac1q$:

    \[ \log\Big(\frac{a^p}{p} + \frac{b^q}{q}\Big) \ge \frac1p\log a^p + \frac1q\log b^q = \log a + \log b = \log(ab). \]

    Exponentiate. For the version with $\varepsilon$, apply the case $p = q = 2$ to $(\sqrt{2\varepsilon}\,a)(b/\sqrt{2\varepsilon})$: $ab \le \frac12(2\varepsilon a^2) + \frac12\frac{b^2}{2\varepsilon} = \varepsilon a^2 + \frac{b^2}{4\varepsilon}$. In general, $ab = \big((\varepsilon p)^{1/p}a\big)\big((\varepsilon p)^{-1/p}b\big) \le \varepsilon a^p + C(\varepsilon)b^q$ with $C(\varepsilon) = \frac{(\varepsilon p)^{-q/p}}{q}$.
  \end{solution}
\end{exercise}

\begin{exercise}
  Prove the interpolation inequality $\norm{f}_{L^r} \le \norm{f}_{L^p}^{\eta}\norm{f}_{L^q}^{1-\eta}$ (for $q < \infty$) from Hölder's inequality.
  \begin{hint}
    Write $|f|^r = |f|^{\eta r}\,|f|^{(1-\eta)r}$ and apply Hölder with exponents $\frac{p}{\eta r}$ and $\frac{q}{(1-\eta)r}$. They are conjugate exactly because $\frac{\eta r}{p} + \frac{(1-\eta)r}{q} = 1$, which is the defining relation for $\eta$.
  \end{hint}
  \begin{solution}
    If $\eta \in \{0,1\}$ then $r = q$ or $r = p$ and there is nothing to prove. Otherwise set $P = \frac{p}{\eta r}$ and $Q = \frac{q}{(1-\eta)r}$. Then $\frac1P + \frac1Q = r\big(\frac\eta p + \frac{1-\eta}{q}\big) = r\cdot\frac1r = 1$, and Hölder gives

    \[ \int|f|^r = \int|f|^{\eta r}|f|^{(1-\eta)r} \le \Big(\int|f|^p\Big)^{\eta r/p}\Big(\int|f|^q\Big)^{(1-\eta)r/q}
             = \norm{f}_{L^p}^{\eta r}\norm{f}_{L^q}^{(1-\eta)r}. \]

    Take the $r$-th root.
  \end{solution}
\end{exercise}

\begin{exercise}
  Let $\rho \ge 0$ with $\int_\Omega \rho = m$, and $\gamma > 2$. Find $\eta$ with $\norm{\rho}_{L^2} \le m^{\eta}\norm{\rho}_{L^\gamma}^{1-\eta}$. Check that $\eta \to 0$ as $\gamma \to 2$ and explain why.
  \begin{hint}
    Solve $\frac12 = \frac{\eta}{1} + \frac{1-\eta}{\gamma}$: $\eta = \frac{\gamma-2}{2(\gamma-1)}$. When $\gamma = 2$ the $L^\gamma$ norm is already the $L^2$ norm, so no help from the mass is needed.
  \end{hint}
  \begin{solution}
    The relation is $\frac12 = \frac\eta1 + \frac{1-\eta}{\gamma}$, i.e. $\eta\big(1 - \frac1\gamma\big) = \frac12 - \frac1\gamma$, so

    \[ \eta = \frac{\gamma-2}{2\gamma}\cdot\frac{\gamma}{\gamma-1} = \frac{\gamma-2}{2(\gamma-1)} \in (0,1). \]

    Since $\rho \ge 0$, $\norm{\rho}_{L^1} = m$, and interpolation gives $\norm{\rho}_{L^2} \le m^\eta\norm{\rho}_{L^\gamma}^{1-\eta}$. As $\gamma \to 2$, $\eta \to 0$: the target space $L^2$ becomes the energy space $L^\gamma$ itself, so no weight on $L^1$ is needed.
  \end{solution}
\end{exercise}

\begin{exercise}
  Absorption. Suppose $X \ge 0$ satisfies $X^2 \le A + BX$ with $A, B \ge 0$. Show $X^2 \le 2A + B^2$. This is the step "move the term to the left-hand side" in every energy estimate.
  \begin{hint}
    By Young, $BX \le \frac12 X^2 + \frac12 B^2$. Substitute and subtract $\frac12 X^2$ from both sides.
  \end{hint}
  \begin{solution}
    By Young, $BX \le \frac12X^2 + \frac12B^2$. Then $X^2 \le A + \frac12X^2 + \frac12B^2$, so $\frac12X^2 \le A + \frac12B^2$ and $X^2 \le 2A + B^2$.
  \end{solution}
\end{exercise}

\subsection*{Where to find a proof}

\begin{itemize}
\item L. C. Evans, \emph{Partial Differential Equations}, 2nd ed., AMS, 2010, Appendix B.2.
\item H. Brezis, \emph{Functional Analysis, Sobolev Spaces and Partial Differential Equations}, Springer, 2011, Chapter 4.
\end{itemize}

\neededfor

% Sobolev embeddings

\section{Sobolev embeddings}\label{p:sobolev-embedding}

\status{stub}
\prereqs{sobolev-spaces}

\subsection*{The idea}

Control of derivatives in $L^p$ gives control of the function itself in a better $L^q$, or even uniformly. This is how an energy bound on $\nabla u$ becomes a bound on $u$ in $L^6$, and how $\rho \in W^{1,q}$ with $q > 3$ gives $\rho \in L^\infty$.

\subsection*{Statement}

\begin{thmbox}{Theorem}
Let $\Omega\subset\R^n$ be bounded with Lipschitz boundary and $1\le p < \infty$.

\begin{itemize}
\item \emph{If $p < n$:} $W^{1,p}(\Omega)\subset L^{q}(\Omega)$ for $1 \le q \le p^* = \frac{np}{n-p}$, with $\norm{u}_{L^{q}} \le C\norm{u}_{W^{1,p}}$.
\item \emph{If $p > n$:} $W^{1,p}(\Omega)\subset C^{0,1-n/p}(\overline\Omega)$; in particular $\norm{u}_{L^\infty}\le C\norm{u}_{W^{1,p}}$.
\item \emph{Zero boundary values:} for $u\in W^{1,p}_0(\Omega)$, $p < n$, the bound holds with $\norm{\nabla u}_{L^p}$ alone on the right. For $n=3$: $\norm{u}_{L^6}\le C\norm{\nabla u}_{L^2}$ for $u\in H^1_0$.
\end{itemize}
\end{thmbox}

In three dimensions, the cases used constantly: $H^1\subset L^6$, $W^{1,q}\subset L^\infty$ for $q> 3$, and $W^{1,\gamma}\subset L^{3\gamma/(3-\gamma)}$ for $\gamma< 3$.

\subsection*{Why it is plausible}

The exponent $p^*$ is forced by scaling. Replace $u(x)$ by $u(\lambda x)$ on $\R^n$. Then $\norm{u}_{L^q}$ scales like $\lambda^{-n/q}$ and $\norm{\nabla u}_{L^p}$ like $\lambda^{1-n/p}$. An inequality $\norm{u}_{L^q}\le C\norm{\nabla u}_{L^p}$ can hold for all $\lambda$ only if $-\frac nq = 1 - \frac np$, that is, $q = p^*$. The same count shows where things break: one derivative in $L^p$ is worth $n/p$ "units of singularity", and when that is less than 1 the function is bounded.

\subsection*{Exercises}

\begin{exercise}
  Carry out the scaling argument. For $u \in C_c^\infty(\R^n)$ and $u_\lambda(x) = u(\lambda x)$, compute $\norm{u_\lambda}_{L^q}$ and $\norm{\nabla u_\lambda}_{L^p}$ in terms of $u$ and $\lambda$, and show that $\norm{u}_{L^q} \le C\norm{\nabla u}_{L^p}$ for all $u$ forces $q = \frac{np}{n-p}$.
  \begin{hint}
    Change variables $y = \lambda x$: $\norm{u_\lambda}_{L^q} = \lambda^{-n/q}\norm{u}_{L^q}$ and $\norm{\nabla u_\lambda}_{L^p} = \lambda^{1-n/p}\norm{\nabla u}_{L^p}$. If the exponents of $\lambda$ differ, let $\lambda \to 0$ or $\lambda \to \infty$ to contradict the inequality.
  \end{hint}
  \begin{solution}
    With $y = \lambda x$, $dx = \lambda^{-n}dy$: $\norm{u_\lambda}_{L^q}^q = \int|u(\lambda x)|^q\,dx = \lambda^{-n}\norm{u}_{L^q}^q$, so $\norm{u_\lambda}_{L^q} = \lambda^{-n/q}\norm{u}_{L^q}$. Since $\nabla u_\lambda(x) = \lambda(\nabla u)(\lambda x)$, similarly $\norm{\nabla u_\lambda}_{L^p} = \lambda^{1-n/p}\norm{\nabla u}_{L^p}$.

    If $\norm{v}_{L^q} \le C\norm{\nabla v}_{L^p}$ for all $v$, apply it to $v = u_\lambda$ with a fixed $u \ne 0$:

    \[ \norm{u}_{L^q} \le C\lambda^{\alpha}\norm{\nabla u}_{L^p}, \qquad \alpha = 1 - \frac np + \frac nq, \quad\text{for every }\lambda > 0. \]

    If $\alpha > 0$, let $\lambda \to 0$; if $\alpha < 0$, let $\lambda \to \infty$; either way $\norm{u}_{L^q} = 0$, a contradiction. So $\alpha = 0$, i.e. $\frac1q = \frac1p - \frac1n$, $q = \frac{np}{n-p}$.
  \end{solution}
\end{exercise}

\begin{exercise}
  The one-dimensional case by hand. For $u \in C^1([0,1])$ with $u(0) = 0$, show $\norm{u}_{L^\infty} \le \norm{u'}_{L^2}$ and $|u(x) - u(y)| \le |x-y|^{1/2}\norm{u'}_{L^2}$. Compare with the theorem for $n = 1$, $p = 2$.
  \begin{hint}
    $u(x) - u(y) = \int_y^x u'$, then Cauchy--Schwarz. Here $p = 2 > n = 1$, and the Hölder exponent $1 - \frac{n}{p} = \frac12$ is exactly what you found.
  \end{hint}
  \begin{solution}
    For $x, y \in [0,1]$, $u(x) - u(y) = \int_y^x u'(t)\,dt$, and Cauchy--Schwarz gives $|u(x) - u(y)| \le |x - y|^{1/2}\big(\int_0^1|u'|^2\big)^{1/2}$. With $y = 0$ and $u(0) = 0$: $|u(x)| \le x^{1/2}\norm{u'}_{L^2} \le \norm{u'}_{L^2}$. For $n = 1$, $p = 2 > n$, the theorem predicts $W^{1,2} \subset C^{0,1-1/2} = C^{0,1/2}$: exactly the exponent $\frac12$ found here.
  \end{solution}
\end{exercise}

\begin{exercise}
  Sharpness in $\R^3$. Using Exercise 2 in \pref{sobolev-spaces}{Sobolev spaces}, show that for every $q > 6$ there is $u \in H^1(B)$ with $u \notin L^q(B)$. So $H^1 \subset L^6$ cannot be improved.
  \begin{hint}
    $|x|^{-s} \in H^1(B)$ for $s < \frac12$, and $|x|^{-s} \notin L^q(B)$ when $sq \ge 3$. Both hold for $s \in [\frac3q, \frac12)$, which is nonempty exactly when $q > 6$.
  \end{hint}
  \begin{solution}
    By Exercise 2 in Sobolev spaces with $n = 3$, $p = 2$: $|x|^{-s} \in H^1(B)$ iff $s < \frac12$. On the other hand $\int_B|x|^{-sq}\,dx = 4\pi\int_0^1 r^{2-sq}\,dr = \infty$ iff $sq \ge 3$. Given $q > 6$, we have $\frac3q < \frac12$, so any $s \in [\frac3q, \frac12)$ gives $u \in H^1(B)\setminus L^q(B)$. For $q \le 6$ no such $s$ exists, consistent with $H^1 \subset L^6$.
  \end{solution}
\end{exercise}

\subsection*{Where to find a proof}

\begin{itemize}
\item L. C. Evans, \emph{Partial Differential Equations}, 2nd ed., AMS, 2010, §5.6 (Gagliardo--Nirenberg--Sobolev inequality and Morrey's inequality).
\item H. Brezis, \emph{Functional Analysis, Sobolev Spaces and Partial Differential Equations}, Springer, 2011, Chapter 9.
\end{itemize}

\neededfor

% Gagliardo–Nirenberg inequalities

\section{Gagliardo--Nirenberg inequalities}\label{p:gagliardo-nirenberg}

\status{stub}
\prereqs{sobolev-embedding, elementary-inequalities}

\subsection*{The idea}

A refinement of Sobolev embedding: an intermediate norm is bounded by a product of a high norm and a low norm, with powers. This is what lets a nonlinear term be split so that the high norm appears with a power below 2 and can be absorbed by Young's inequality.

\subsection*{Statement}

\begin{thmbox}{Theorem (on $\R^n$)}
Let $1\le q, r\le\infty$, $j < m$ integers, and $\frac{j}{m}\le a\le 1$ with

\[ \frac1p = \frac jn + a\Big(\frac1r - \frac mn\Big) + \frac{1-a}{q}. \]

Then $\norm{\nabla^j u}_{L^p}\le C\norm{\nabla^m u}_{L^r}^{a}\norm{u}_{L^q}^{1-a}$, except that $a = 1$ is not allowed when $1< r< \infty$ and $m - j - n/r$ is a nonnegative integer. On a bounded Lipschitz domain the same holds with an extra lower-order term, e.g. $+\,C\norm{u}_{L^q}$, on the right.
\end{thmbox}

\begin{exbox}{The form used most in 3D}
For $f\in H^1(\Omega)$, $\Omega\subset\R^3$ bounded and Lipschitz: $\norm{f}_{L^3}\le C\norm{f}_{L^2}^{1/2}\norm{f}_{H^1}^{1/2}$. Applied to $f=\nabla u$: $\norm{\nabla u}_{L^3}\le C\norm{\nabla u}_{L^2}^{1/2}\norm{\nabla u}_{H^1}^{1/2}$.
\end{exbox}

\subsection*{Why it is plausible}

The example is just two facts in a row: interpolate $L^3$ between $L^2$ and $L^6$ (the weights are $\frac12,\frac12$ because $\frac13 = \frac{1/2}{2} + \frac{1/2}{6}$), then bound $L^6$ by $H^1$ with Sobolev. The general exponent relation is again forced by scaling, exactly as for Sobolev embedding.

\subsection*{Exercises}

\begin{exercise}
  In $\R^3$, prove $\norm{f}_{L^4} \le C\norm{f}_{L^2}^{1/4}\norm{f}_{H^1}^{3/4}$ by interpolation and Sobolev. Then check the exponent $a = \frac34$ against the general formula with $j = 0$, $m = 1$, $r = q = 2$, $n = 3$.
  \begin{hint}
    Interpolate $L^4$ between $L^2$ and $L^6$: $\frac14 = \frac{\eta}{2} + \frac{1-\eta}{6}$ gives $\eta = \frac14$. Then $\norm{f}_{L^6} \le C\norm{f}_{H^1}$. The formula reads $\frac14 = a(\frac12 - \frac13) + \frac{1-a}{2} = \frac12 - \frac a3$.
  \end{hint}
  \begin{solution}
    Interpolation: $\frac14 = \frac\eta2 + \frac{1-\eta}{6}$; multiplying by $12$, $3 = 6\eta + 2 - 2\eta$, so $\eta = \frac14$ and $\norm{f}_{L^4} \le \norm{f}_{L^2}^{1/4}\norm{f}_{L^6}^{3/4}$. Sobolev on a bounded Lipschitz domain in $\R^3$ gives $\norm{f}_{L^6} \le C\norm{f}_{H^1}$, so $\norm{f}_{L^4} \le C\norm{f}_{L^2}^{1/4}\norm{f}_{H^1}^{3/4}$.

    The general formula with $j = 0$, $m = 1$, $r = q = 2$, $n = 3$ reads $\frac14 = a\big(\frac12 - \frac13\big) + \frac{1-a}{2} = \frac12 - \frac a3$, so $a = \frac34$, the same exponent.
  \end{solution}
\end{exercise}

\begin{exercise}
  Find the exponent $a$ in $\norm{f}_{L^4(\R^2)} \le C\norm{\nabla f}_{L^2}^{a}\norm{f}_{L^2}^{1-a}$ (Ladyzhenskaya's inequality) from the formula, and confirm it by the scaling $f(x) \mapsto f(\lambda x)$.
  \begin{hint}
    With $n = 2$ the formula gives $\frac14 = a\cdot 0 + \frac{1-a}{2}$, so $a = \frac12$. Under scaling, $\norm{f}_{L^4} \sim \lambda^{-1/2}$, $\norm{\nabla f}_{L^2} \sim \lambda^0$ and $\norm{f}_{L^2} \sim \lambda^{-1}$.
  \end{hint}
  \begin{solution}
    For $n = 2$: $\frac14 = a\big(\frac12 - \frac12\big) + \frac{1-a}{2}$, so $a = \frac12$. Scaling $f_\lambda(x) = f(\lambda x)$ on $\R^2$: $\norm{f_\lambda}_{L^4} = \lambda^{-1/2}\norm{f}_{L^4}$, $\norm{\nabla f_\lambda}_{L^2} = \lambda^{0}\norm{\nabla f}_{L^2}$, $\norm{f_\lambda}_{L^2} = \lambda^{-1}\norm{f}_{L^2}$. For the inequality to survive all $\lambda$, the powers must match: $-\frac12 = 0\cdot a - (1-a)$, again $a = \frac12$.
  \end{solution}
\end{exercise}

\begin{exercise}
  The convection term in Step 1 of the case study. For $u \in H^1_0 \cap H^2(\Omega)$, $\Omega \subset \R^3$, show $\norm{u\cdot\nabla u}_{L^2} \le C\norm{\nabla u}_{L^2}^{3/2}\norm{\nabla u}_{H^1}^{1/2}$. Then show that $\norm{\nabla u}_{H^1} \le C\big(A + \norm{u\cdot\nabla u}_{L^2}\big)$ implies $\norm{\nabla u}_{H^1} \le C'\big(A + \norm{\nabla u}_{L^2}^3\big)$.
  \begin{hint}
    Hölder: $\norm{u\cdot\nabla u}_{L^2} \le \norm{u}_{L^6}\norm{\nabla u}_{L^3}$. Use $\norm{u}_{L^6} \le C\norm{\nabla u}_{L^2}$ and the 3D form above for $\nabla u$. For the second part, write $X = \norm{\nabla u}_{H^1}$, $K = \norm{\nabla u}_{L^2}$, and use Young: $CK^{3/2}X^{1/2} \le \frac12 X + C''K^3$. The power $\frac12$ of $X$ is below $1$: this is "absorbed".
  \end{hint}
  \begin{solution}
    Hölder with $\frac12 = \frac16 + \frac13$: $\norm{u\cdot\nabla u}_{L^2} \le \norm{u}_{L^6}\norm{\nabla u}_{L^3}$. Since $u \in H^1_0$, $\norm{u}_{L^6} \le C\norm{\nabla u}_{L^2}$. Since $u \in H^2$, $\nabla u \in H^1$, and the 3D form gives $\norm{\nabla u}_{L^3} \le C\norm{\nabla u}_{L^2}^{1/2}\norm{\nabla u}_{H^1}^{1/2}$. Multiplying, $\norm{u\cdot\nabla u}_{L^2} \le C_1K^{3/2}X^{1/2}$ with $K = \norm{\nabla u}_{L^2}$, $X = \norm{\nabla u}_{H^1}$.

    Then $X \le CA + CC_1K^{3/2}X^{1/2}$. By Young with $p = q = 2$, $CC_1K^{3/2}X^{1/2} \le \frac12X + \frac12(CC_1)^2K^3$. So $\frac12X \le CA + \frac12(CC_1)^2K^3$, i.e. $X \le C'(A + K^3)$. The term was absorbed because $X$ appeared on the right with a power below $1$.
  \end{solution}
\end{exercise}

\subsection*{Where to find a proof}

\begin{itemize}
\item L. Nirenberg, On elliptic partial differential equations, \emph{Ann. Scuola Norm. Sup. Pisa} 13 (1959) 115--162 (the original statement).
\item A. Fiorenza, M. R. Formica, T. Roskovec, F. Soudský, Detailed proof of classical Gagliardo--Nirenberg interpolation inequality with historical remarks, \emph{Z. Anal. Anwend.} 40 (2021) 217--236.
\end{itemize}

\neededfor

% Gronwall’s inequality and continuity arguments
% Converted from notes/pde/gronwall.html. Edit freely: this file is now the master copy.

\section{Gronwall's inequality and continuity arguments}\label{p:gronwall}

\subsection*{The idea}

Energy estimates usually end with a differential inequality: the quantity you control grows at a rate bounded by itself. Gronwall's inequality turns this into an explicit bound. Its nonlinear relative, the continuity argument, gives bounds for a short time when the growth rate is superlinear.

\subsection*{Statement}

\begin{thmbox}{Theorem (Gronwall, differential form)}
Let $\eta\ge0$ be absolutely continuous on $[0,T]$ and $\phi,\psi\ge0$ integrable, with $\eta'(t)\le\phi(t)\eta(t)+\psi(t)$ for a.e. $t$. Then

\[ \eta(t)\le e^{\int_0^t\phi}\Big(\eta(0)+\int_0^t\psi\Big)\quad\text{for all } t\in[0,T]. \]

In particular, if $\psi = 0$ and $\eta(0)=0$, then $\eta\equiv0$.
\end{thmbox}

\begin{thmbox}{Continuity argument}
Let $\Phi\ge0$ be continuous and nondecreasing on $[0,T]$, and $M$ nondecreasing, with $\Phi(t)\le A\exp\big(C\int_0^t M(\Phi(s))\,ds\big)$ for all $t$. If $0< T_0< \frac{\ln 2}{C\,M(2A)}$, then $\Phi(t)< 2A$ on $[0,\min(T_0,T)]$.
\end{thmbox}

\subsection*{Why it is plausible}

For Gronwall: $\big(e^{-\int_0^t\phi}\eta\big)'\le e^{-\int_0^t\phi}\psi$, then integrate. For the continuity argument: as long as $\Phi\le 2A$, the right-hand side is at most $A\exp(C\,t\,M(2A))< 2A$ for $t\le T_0$. So $\Phi$ can never reach the value $2A$ before time $T_0$: the first time it did, the inequality would give a strict contradiction. No formula for $M$ is needed, which is why papers write "for some $M$" and let it change from line to line.

\subsection*{Exercises}

\begin{exercise}
  Prove the differential form of Gronwall's inequality.
  \begin{hint}
    Let $\Phi(t) = \int_0^t \phi$. Then $\big(e^{-\Phi}\eta\big)' = e^{-\Phi}(\eta' - \phi\eta) \le e^{-\Phi}\psi \le \psi$ a.e. Integrate from $0$ to $t$ and multiply by $e^{\Phi(t)}$.
  \end{hint}
  \begin{solution}
    Let $\Phi(t) = \int_0^t\phi$, absolutely continuous with $\Phi' = \phi$ a.e. The product $e^{-\Phi}\eta$ is absolutely continuous and, a.e.,

    \[ \big(e^{-\Phi}\eta\big)' = e^{-\Phi}\big(\eta' - \phi\eta\big) \le e^{-\Phi}\psi \le \psi, \]

    using $\Phi \ge 0$ and $\psi \ge 0$. Integrating from $0$ to $t$: $e^{-\Phi(t)}\eta(t) \le \eta(0) + \int_0^t\psi$. Multiply by $e^{\Phi(t)}$. If $\psi = 0$ and $\eta(0) = 0$ the right side is $0$, and $\eta \ge 0$ forces $\eta \equiv 0$.
  \end{solution}
\end{exercise}

\begin{exercise}
  Compare the continuity argument with an exact solution. Let $y' = y^2$, $y(0) = A > 0$. (a) Solve it and show $y(t) = A\exp\big(\int_0^t y\big)$, so the continuity argument applies with $C = 1$, $M(s) = s$. (b) Which time does the argument guarantee, and when does $y$ actually reach $2A$?
  \begin{hint}
    (a) $y = \frac{A}{1 - At}$, and $\int_0^t y = -\log(1 - At)$. (b) The argument gives $y < 2A$ for $t < \frac{\ln 2}{2A} \approx \frac{0.35}{A}$. In fact $y = 2A$ at $t = \frac{1}{2A}$, and $y$ blows up at $t = \frac1A$. The argument is not sharp, but it needs no formula for the solution.
  \end{hint}
  \begin{solution}
    (a) Separating variables, $-\frac1y = t - \frac1A$, so $y(t) = \frac{A}{1 - At}$ for $t < \frac1A$. Then $\int_0^t y = \int_0^t\frac{A}{1-As}\,ds = -\log(1 - At)$, and $A\exp\big(\int_0^ty\big) = \frac{A}{1-At} = y(t)$. So $\Phi = y$ satisfies the hypothesis of the continuity argument with $C = 1$, $M(s) = s$ (with equality).

    (b) The argument gives $y < 2A$ on $[0, T_0]$ for any $T_0 < \frac{\ln 2}{M(2A)} = \frac{\ln 2}{2A} \approx \frac{0.347}{A}$. Exactly, $\frac{A}{1-At} = 2A$ at $t = \frac{1}{2A}$, and $y$ blows up at $t = \frac1A$. The guaranteed time is about $70\%$ of the true one, without solving anything.
  \end{solution}
\end{exercise}

\begin{exercise}
  Why $K \in L^1$ matters. Show that $y(t) = t^2$ and $y = 0$ both solve $y' = 2\sqrt{y}$, $y(0) = 0$. Writing the equation as $y' = K(t)y$, explain why Gronwall cannot rule out the nonzero solution.
  \begin{hint}
    For $y = t^2$, $K(t) = \frac{2\sqrt{y}}{y} = \frac{2}{t}$, which is not integrable near $t = 0$.
  \end{hint}
  \begin{solution}
    For $t \ge 0$, $y = t^2$ gives $y' = 2t = 2\sqrt{t^2} = 2\sqrt y$, and $y(0) = 0$; $y = 0$ is also a solution. For the nonzero solution, $y' = K(t)y$ with $K = \frac{2\sqrt y}{y} = \frac2t$, and $\int_0^T\frac2t\,dt = \infty$. Gronwall needs $\phi = K \in L^1(0,T)$: its bound $e^{\int_0^tK}\cdot y(0)$ becomes the meaningless $\infty\cdot 0$. The non-Lipschitz right side $2\sqrt y$ at $y = 0$ is exactly what lets a solution leave $0$.
  \end{solution}
\end{exercise}

\subsection*{Where to find a proof}

\begin{itemize}
\item L. C. Evans, \emph{Partial Differential Equations}, 2nd ed., AMS, 2010, Appendix B.2 (Gronwall's inequality, differential and integral forms).
\item T. Tao, \emph{Nonlinear Dispersive Equations: Local and Global Analysis}, AMS, 2006, §1.3 (continuity arguments).
\end{itemize}

\neededfor{Energy estimates and a priori bounds~(\S\ref{p:energy-estimates}), Uniqueness by difference estimates~(\S\ref{p:uniqueness-energy})}


\chapter{Linear building blocks}
% Elliptic regularity: the Lamé system and the Neumann Laplacian

\section{Elliptic regularity: the Lamé system and the Neumann Laplacian}\label{p:elliptic-regularity}

\status{stub}
\prereqs{sobolev-spaces}

\subsection*{The idea}

For an elliptic operator, knowing the operator applied to $u$ in $L^p$ gives all second derivatives of $u$ in $L^p$. In evolution problems this is used "frozen in time": rewrite the equation as (elliptic operator)$(u)$ = (everything else), and read off two derivatives of $u$ from bounds on everything else.

\subsection*{Statement}

\begin{thmbox}{Theorem (Lamé system, Dirichlet condition)}
\index{Lamé system}\index{elliptic regularity!Lamé system}%
Let $\Omega\subset\R^3$ be bounded with smooth boundary, $\Lame u = \mu\Delta u+(\lambda+\mu)\nabla\dv u$ with $\mu> 0$ and $\lambda+2\mu> 0$, and $1< p< \infty$. If $u\in W^{1,p}_0(\Omega)$ and $\Lame u = F\in L^p$, then $u\in W^{2,p}(\Omega)$ and $\norm{u}_{W^{2,p}}\le C\norm{F}_{L^p}$, with $C$ depending on $\Omega, p, \mu, \lambda$.
\end{thmbox}

\begin{thmbox}{Theorem (Laplacian, Neumann condition)}
\index{elliptic regularity!Neumann Laplacian}%
If $\phi\in H^1(\Omega)$ is a weak solution of $-\Delta\phi = f\in L^2$ with $\partial_\nu\phi = 0$ on $\partial\Omega$, then $\phi\in H^2(\Omega)$ and $\norm{\phi}_{H^2}\le C(\norm{f}_{L^2}+\norm{\phi}_{L^2})$. The same holds in $W^{2,p}$ for $1< p< \infty$.
\end{thmbox}

\subsection*{Why it is plausible}

On all of $\R^n$ and for the Laplacian in $L^2$ this is integration by parts twice: $\int|\nabla^2 u|^2 = \int|\Delta u|^2$. The condition $\lambda+2\mu> 0$ is what makes the same identity work for the Lamé operator: it splits into $\mu\Delta$ on the divergence-free part of $u$ and $(\lambda+2\mu)\Delta$ on the gradient part, and both coefficients must be positive. In $L^p$ with $p\ne2$ the identity is replaced by Calderón--Zygmund theory, and near the boundary by flattening it; this is the hard part of the proof.

\subsection*{Exercises}

\begin{exercise}
  For $u \in C_c^\infty(\R^n)$, prove $\int|\nabla^2 u|^2\,dx = \int|\Delta u|^2\,dx$.
  \begin{hint}
    Integrate by parts twice: $\int \partial_i\partial_j u\,\partial_i\partial_j u = -\int \partial_j u\,\partial_i\partial_i\partial_j u = \int \partial_j\partial_j u\,\partial_i\partial_i u$, summed over $i, j$.
  \end{hint}
  \begin{solution}
    Sum over $i, j$ and integrate by parts first in $x_i$, then in $x_j$; there are no boundary terms because $u$ has compact support:

    \[ \int|\nabla^2u|^2 = \sum_{i,j}\int\partial_i\partial_ju\,\partial_i\partial_ju = -\sum_{i,j}\int\partial_ju\,\partial_i\partial_i\partial_ju
             = \sum_{i,j}\int\partial_j\partial_ju\,\partial_i\partial_iu = \int(\Delta u)^2. \]
  \end{solution}
\end{exercise}

\begin{exercise}
  Where $\lambda + 2\mu > 0$ comes from. For $u \in C_c^\infty(\R^3;\R^3)$, show $-\int \Lame u\cdot u = \mu\int|\nabla\times u|^2 + (\lambda + 2\mu)\int(\dv u)^2$, using $\int|\nabla u|^2 = \int|\nabla\times u|^2 + \int(\dv u)^2$. Then explain why the right-hand side controls $\int|\nabla u|^2$ exactly when $\mu > 0$ and $\lambda + 2\mu > 0$.
  \begin{hint}
    Integrate by parts: $-\int \Lame u\cdot u = \mu\int|\nabla u|^2 + (\lambda+\mu)\int(\dv u)^2$, then substitute the identity. Test with a curl-free field ($u = \nabla\psi$) and with a divergence-free field to see that each coefficient must be positive.
  \end{hint}
  \begin{solution}
    Integrating by parts, $-\int\Delta u\cdot u = \int|\nabla u|^2$ and $-\int\nabla\dv u\cdot u = \int(\dv u)^2$, so

    \[ -\int\Lame u\cdot u = \mu\int|\nabla u|^2 + (\lambda+\mu)\int(\dv u)^2. \]

    For compactly supported $u$, $-\Delta u = \nabla\times\nabla\times u - \nabla\dv u$; dotting with $u$ and integrating by parts gives $\int|\nabla u|^2 = \int|\nabla\times u|^2 + \int(\dv u)^2$. Substituting,

    \[ -\int\Lame u\cdot u = \mu\int|\nabla\times u|^2 + (\lambda+2\mu)\int(\dv u)^2. \]

    If $\mu > 0$ and $\lambda + 2\mu > 0$, this is at least $\min(\mu,\lambda+2\mu)\int|\nabla u|^2$. Conversely, take $u = \nabla\psi$ with $\psi \in C_c^\infty$: then $\nabla\times u = 0$ and the form equals $(\lambda+2\mu)\int(\Delta\psi)^2$, so $\lambda + 2\mu > 0$ is needed. Take $u = \nabla\times\Psi$: then $\dv u = 0$ and the form equals $\mu\int|\nabla\times u|^2$, so $\mu > 0$ is needed.
  \end{solution}
\end{exercise}

\begin{exercise}
  The Neumann problem. Show that constants solve $-\Delta\phi = 0$ with $\partial_\nu\phi = 0$, so the estimate cannot hold without the term $\norm{\phi}_{L^2}$. Show also that $-\Delta\phi = f$ with the Neumann condition can have a solution only if $\int_\Omega f = 0$.
  \begin{hint}
    Integrate the equation over $\Omega$ and use the divergence theorem: $\int_\Omega \Delta\phi = \int_{\partial\Omega}\partial_\nu\phi = 0$.
  \end{hint}
  \begin{solution}
    A constant $\phi = c$ has $\Delta\phi = 0$ and $\partial_\nu\phi = 0$. An estimate $\norm{\phi}_{H^2} \le C\norm{f}_{L^2}$ would give $\norm{c}_{H^2} \le 0$, false for $c \ne 0$; the term $\norm{\phi}_{L^2}$ is what accounts for constants.

    If $-\Delta\phi = f$ with $\partial_\nu\phi = 0$, the divergence theorem gives $\int_\Omega f = -\int_\Omega\Delta\phi = -\int_{\partial\Omega}\partial_\nu\phi = 0$. So $\int f = 0$ is necessary. (It is also sufficient, and the solution is unique up to a constant.)
  \end{solution}
\end{exercise}

\subsection*{Where to find a proof}

\begin{itemize}
\item L. C. Evans, \emph{Partial Differential Equations}, 2nd ed., AMS, 2010, §6.3 (H\textsuperscript{2} regularity for scalar equations, Dirichlet case).
\item S. Agmon, A. Douglis, L. Nirenberg, Estimates near the boundary for solutions of elliptic partial differential equations satisfying general boundary conditions II, \emph{Comm. Pure Appl. Math.} 17 (1964) 35--92 (systems, including Lamé).
\item D. Gilbarg, N. S. Trudinger, \emph{Elliptic Partial Differential Equations of Second Order}, Springer, Chapter 9 ($W^{2,p}$ estimates).
\end{itemize}

\neededfor

% Maximal regularity for parabolic equations

\section{Maximal regularity for parabolic equations}\label{p:parabolic-regularity}

\status{stub}
\prereqs{bochner-spaces, elliptic-regularity}

\subsection*{The idea}

For a heat-type equation $\phi_t - D\Delta\phi = f$, a source $f$ in $L^p$ in space and time gives both $\phi_t$ and all second space derivatives of $\phi$ in the same $L^p$. "Maximal" means nothing is lost: the solution is exactly as regular as the equation allows.

\subsection*{Statement}

\begin{thmbox}{Theorem}
Let $\Omega\subset\R^3$ be bounded with smooth boundary, $1< p< \infty$, $\sigma, D> 0$, $b\ge0$, and $Q_T=\Omega\times(0,T)$. Consider

\[ \sigma\phi_t - D\Delta\phi + b\phi = f \text{ in } Q_T,\qquad \partial_\nu\phi = 0 \text{ on } \partial\Omega,\qquad \phi(0)=\phi_0. \]

If $f\in L^p(Q_T)$ and $\phi_0\in W^{2-2/p,\,p}(\Omega)$ (with the compatible boundary condition when $p> 3$), then there is a unique solution with $\norm{\phi_t}_{L^p(Q_T)}+\norm{\nabla^2\phi}_{L^p(Q_T)}\le C\big(\norm{f}_{L^p(Q_T)}+\norm{\phi_0}_{W^{2-2/p,p}}\big)$.
\end{thmbox}

\subsection*{Why it is plausible}

The equation balances one time derivative against two space derivatives, so whatever regularity $f$ has, $\phi_t$ and $\nabla^2\phi$ share it. The odd exponent $2-\frac2p$ for the initial data is the trace: a function with $\phi_t,\nabla^2\phi\in L^p$ has, at a fixed time, exactly $2-\frac2p$ space derivatives in $L^p$ (one loses $\frac1p$ of a time derivative, which is worth $\frac2p$ space derivatives).

\subsection*{Exercises}

\begin{exercise}
  The case $p = 2$ by hand. Let $\phi_t - \Delta\phi = f$ in $\Omega\times(0,T)$ with $\partial_\nu\phi = 0$ and $\phi(0) = \phi_0$. Test with $-\Delta\phi$ to show that for every $t$, $\norm{\nabla\phi(t)}_{L^2}^2 + \int_0^t\norm{\Delta\phi}_{L^2}^2 \le \norm{\nabla\phi_0}_{L^2}^2 + \int_0^t\norm{f}_{L^2}^2$, then bound $\int_0^T\norm{\phi_t}_{L^2}^2$. The data needed is $\phi_0 \in H^1 = W^{2-2/p,\,p}$ for $p = 2$.
  \begin{hint}
    $\int\phi_t(-\Delta\phi) = \int\nabla\phi_t\cdot\nabla\phi = \frac12\frac{d}{dt}\norm{\nabla\phi}_{L^2}^2$, with no boundary term because $\partial_\nu\phi = 0$. Bound $-\int f\Delta\phi \le \frac12\norm{f}^2 + \frac12\norm{\Delta\phi}^2$ and absorb. For $\phi_t$, use the equation: $\phi_t = f + \Delta\phi$. The $H^2$ bound then follows from elliptic regularity.
  \end{hint}
  \begin{solution}
    For a smooth solution, multiply the equation by $-\Delta\phi$ and integrate. Since $\partial_\nu\phi = 0$,

    \[ \int\phi_t(-\Delta\phi) = \int\nabla\phi_t\cdot\nabla\phi = \frac12\frac{d}{dt}\norm{\nabla\phi}_{L^2}^2, \qquad
             \int(-\Delta\phi)(-\Delta\phi) = \norm{\Delta\phi}_{L^2}^2, \]

    and $-\int f\Delta\phi \le \frac12\norm{f}_{L^2}^2 + \frac12\norm{\Delta\phi}_{L^2}^2$. Absorbing, $\frac{d}{dt}\norm{\nabla\phi}_{L^2}^2 + \norm{\Delta\phi}_{L^2}^2 \le \norm{f}_{L^2}^2$, and integrating from $0$ to $t$:

    \[ \norm{\nabla\phi(t)}_{L^2}^2 + \int_0^t\norm{\Delta\phi}_{L^2}^2 \le \norm{\nabla\phi_0}_{L^2}^2 + \int_0^t\norm{f}_{L^2}^2. \]

    From the equation, $\phi_t = f + \Delta\phi$, so $\norm{\phi_t}^2 \le 2\norm{f}^2 + 2\norm{\Delta\phi}^2$ and $\int_0^T\norm{\phi_t}_{L^2}^2 \le 2\norm{\nabla\phi_0}_{L^2}^2 + 4\int_0^T\norm{f}_{L^2}^2$. Elliptic regularity for the Neumann Laplacian turns $\Delta\phi \in L^2$ into $\phi \in H^2$ (with $\norm{\phi}_{L^2}$, bounded by testing with $\phi$). The data enter only through $\nabla\phi_0 \in L^2$: $\phi_0 \in H^1 = W^{2-2/p,p}$ for $p = 2$.
  \end{solution}
\end{exercise}

\begin{exercise}
  Show that the heat equation is invariant under $x \mapsto \lambda x$, $t \mapsto \lambda^2 t$, and that $\norm{\phi_t}_{L^p(Q)}$ and $\norm{\nabla^2\phi}_{L^p(Q)}$ scale with the same power of $\lambda$. This is why one time derivative is worth two space derivatives.
  \begin{hint}
    For $\phi_\lambda(x,t) = \phi(\lambda x, \lambda^2 t)$, both $\partial_t\phi_\lambda$ and $\nabla^2\phi_\lambda$ carry a factor $\lambda^2$, and the change of variables in $Q$ contributes $\lambda^{-(n+2)/p}$ to both norms.
  \end{hint}
  \begin{solution}
    Let $\phi_\lambda(x,t) = \phi(\lambda x, \lambda^2t)$. By the chain rule, $\partial_t\phi_\lambda = \lambda^2(\phi_t)(\lambda x,\lambda^2t)$ and $\nabla^2\phi_\lambda = \lambda^2(\nabla^2\phi)(\lambda x,\lambda^2t)$; so $\phi_t = \Delta\phi$ implies the same for $\phi_\lambda$. Changing variables $(y,s) = (\lambda x,\lambda^2t)$, $dx\,dt = \lambda^{-n-2}dy\,ds$:

    \[ \norm{\partial_t\phi_\lambda}_{L^p} = \lambda^{2 - \frac{n+2}{p}}\norm{\phi_t}_{L^p}, \qquad
             \norm{\nabla^2\phi_\lambda}_{L^p} = \lambda^{2 - \frac{n+2}{p}}\norm{\nabla^2\phi}_{L^p}, \]

    the same power, while $\norm{\nabla\phi_\lambda}_{L^p}$ carries $\lambda^{1 - \frac{n+2}{p}}$. One $t$-derivative costs what two $x$-derivatives cost.
  \end{solution}
\end{exercise}

\subsection*{Where to find a proof}

\begin{itemize}
\item E. Feireisl, A. Novotný, \emph{Singular Limits in Thermodynamics of Viscous Fluids}, Birkhäuser, 2009, Theorem 10.22.
\item A. Novotný, I. Straškraba, \emph{Introduction to the Mathematical Theory of Compressible Flow}, Oxford, 2004, Lemma 7.37.
\end{itemize}

\neededfor

% The transport (continuity) equation

\section{The transport (continuity) equation}\label{p:transport-equation}

\status{stub}
\prereqs{sobolev-spaces, sobolev-embedding}

\subsection*{The idea}

The continuity equation $\rho_t+\dv(\rho u)=0$ carries the density along the flow of $u$. Unlike parabolic equations it never smooths: $\rho(t)$ is exactly as regular as $\rho_0$ and the flow allow. All information about $\rho$ comes from following particle paths, and the cost is paid in $\nabla u$.

\subsection*{Statement}

\begin{thmbox}{Properties}
Let $u$ be smooth with $u=0$ on $\partial\Omega$ and $\rho$ a smooth solution of $\rho_t+\dv(\rho u)=0$. Then:

\begin{itemize}
\item \emph{Mass:} $\int_\Omega\rho(t)\,dx=\int_\Omega\rho_0\,dx$.
\item \emph{Sign and lower bound:} if $\rho_0\ge\delta> 0$ then $\rho(x,t)\ge\delta\exp\big(-\int_0^t\norm{\dv u}_{L^\infty}ds\big)> 0$; if $\rho_0\ge0$ then $\rho\ge0$, and vacuum stays vacuum.
\item \emph{Derivatives:} for $q\in(3,6]$, $\norm{\rho(t)}_{H^1\cap W^{1,q}}\le\norm{\rho_0}_{H^1\cap W^{1,q}}\exp\big(C\int_0^t\norm{\nabla u}_{H^1\cap W^{1,q}}ds\big)$.
\end{itemize}
\end{thmbox}

\subsection*{Why it is plausible}

Along a particle path $\dot X = u(X,t)$, the equation reads $\frac{d}{dt}\rho(X(t),t) = -\rho\,\dv u$, so $\rho$ is multiplied by $\exp(-\int\dv u)$: this gives positivity and the lower bound. Differentiating the equation, $\nabla\rho$ satisfies a transport equation with source $-\nabla u\cdot\nabla\rho - \rho\nabla\dv u$; the first term needs $\nabla u\in L^\infty$, which is why one asks for $\nabla u\in W^{1,q}$ with $q> 3$ (Sobolev embedding), and the second needs $\nabla^2 u\in L^q$.

\subsection*{Exercises}

\begin{exercise}
  Mass and positivity. Show $\frac{d}{dt}\int_\Omega\rho = 0$ when $u = 0$ on $\partial\Omega$. Along a path $\dot X = u(X,t)$, show $\frac{d}{dt}\rho(X(t),t) = -\rho\,\dv u$ and solve this to get the lower bound in the statement.
  \begin{hint}
    $\int\dv(\rho u) = \int_{\partial\Omega}\rho\,u\cdot\nu = 0$. For the path, expand $\dv(\rho u) = u\cdot\nabla\rho + \rho\,\dv u$ and use the chain rule; then $\rho(X(t),t) = \rho_0(X(0))\exp\big(-\int_0^t \dv u(X(s),s)\,ds\big)$.
  \end{hint}
  \begin{solution}
    $\frac{d}{dt}\int_\Omega\rho = -\int_\Omega\dv(\rho u) = -\int_{\partial\Omega}\rho\,u\cdot\nu = 0$ since $u = 0$ on $\partial\Omega$.

    Since $\dv(\rho u) = u\cdot\nabla\rho + \rho\,\dv u$, along $\dot X = u(X,t)$:

    \[ \frac{d}{dt}\rho(X(t),t) = \rho_t + \nabla\rho\cdot\dot X = \rho_t + u\cdot\nabla\rho = -\rho\,\dv u. \]

    This linear ODE gives $\rho(X(t),t) = \rho_0(X(0))\exp\big(-\int_0^t\dv u(X(s),s)\,ds\big)$. Because $u$ is smooth and vanishes on $\partial\Omega$, every point $(x,t)$ lies on exactly one path starting in $\overline\Omega$. Hence $\rho_0 \ge \delta$ gives $\rho(x,t) \ge \delta\exp\big(-\int_0^t\norm{\dv u}_{L^\infty}\big)$; $\rho_0 \ge 0$ gives $\rho \ge 0$; and $\rho_0 = 0$ at the start of a path gives $\rho = 0$ along it.
  \end{solution}
\end{exercise}

\begin{exercise}
  An explicit compressing flow. On $\R$ take $u(x) = -x$. Check that $\rho(x,t) = e^t\rho_0(xe^t)$ solves $\rho_t + (\rho u)_x = 0$. Show that the mass is conserved, that $\norm{\rho_x(t)}_{L^\infty} = e^{2t}\norm{\rho_0'}_{L^\infty}$, and that a jump in $\rho_0$ is never smoothed out.
  \begin{hint}
    Differentiate directly. The paths are $X(t) = x_0e^{-t}$: particles are squeezed together, the density grows like $e^t$, and its derivative picks up one more factor $e^t$ from the squeezing. A jump at $x_0$ moves to $x_0e^{-t}$ and stays a jump.
  \end{hint}
  \begin{solution}
    With $\rho = e^t\rho_0(xe^t)$: $\rho_t = e^t\rho_0(xe^t) + xe^{2t}\rho_0'(xe^t)$ and $(\rho u)_x = (-x\rho)_x = -\rho - x\rho_x = -e^t\rho_0(xe^t) - xe^{2t}\rho_0'(xe^t)$. The sum is $0$.

    Mass: $\int e^t\rho_0(xe^t)\,dx = \int\rho_0(y)\,dy$. Derivative: $\rho_x = e^{2t}\rho_0'(xe^t)$, so $\norm{\rho_x(t)}_{L^\infty} = e^{2t}\norm{\rho_0'}_{L^\infty}$. Here $\nabla u = -1$, so the growth is $\exp\big(2\int_0^t\norm{\nabla u}_{L^\infty}\big)$, of the form in the statement. If $\rho_0$ jumps at $x_0$, then $\rho(\cdot,t)$ jumps at $x_0e^{-t}$ by $e^t$ times the original jump: the equation moves and amplifies the jump but never smooths it.
  \end{solution}
\end{exercise}

\begin{exercise}
  Differentiate $\rho_t + \dv(\rho u) = 0$ in $x_i$ to get $(\partial_i\rho)_t + \dv(\partial_i\rho\,u) = -\partial_i u\cdot\nabla\rho - \rho\,\partial_i\dv u$. Multiply by $|\nabla\rho|^{q-2}\partial_i\rho$, sum over $i$, and show $\frac{d}{dt}\norm{\nabla\rho}_{L^q} \le C\norm{\nabla u}_{L^\infty}\norm{\nabla\rho}_{L^q} + \norm{\rho}_{L^\infty}\norm{\nabla^2u}_{L^q}$. Why does this need $q > 3$ in the case study?
  \begin{hint}
    The term $\int\dv(\partial_i\rho\,u)|\nabla\rho|^{q-2}\partial_i\rho$ integrates by parts to a multiple of $\int\dv u\,|\nabla\rho|^q$. To close with Gronwall you need $\nabla u \in L^1(0,T;L^\infty)$; in 3D, $W^{1,q} \subset L^\infty$ requires $q > 3$.
  \end{hint}
  \begin{solution}
    Differentiating, $\partial_i\dv(\rho u) = \dv(\partial_i\rho\,u) + \partial_iu\cdot\nabla\rho + \rho\,\partial_i\dv u$, which gives the stated equation. Multiply it by $|\nabla\rho|^{q-2}\partial_i\rho$ and sum over $i$. The time term is $\frac1q\partial_t|\nabla\rho|^q$. Since $\dv(\partial_i\rho\,u) = u\cdot\nabla\partial_i\rho + \partial_i\rho\,\dv u$, the transport term is $\frac1qu\cdot\nabla|\nabla\rho|^q + \dv u\,|\nabla\rho|^q$, whose integral is $(1 - \frac1q)\int\dv u\,|\nabla\rho|^q$ (integrate the first piece by parts; $u = 0$ on $\partial\Omega$). Hence, with Hölder on the last term,

    \[ \frac1q\frac{d}{dt}\norm{\nabla\rho}_{L^q}^q \le C\norm{\nabla u}_{L^\infty}\norm{\nabla\rho}_{L^q}^q
             + \norm{\rho}_{L^\infty}\norm{\nabla^2u}_{L^q}\norm{\nabla\rho}_{L^q}^{q-1}. \]

    Since $\frac1q\frac{d}{dt}N^q = N^{q-1}N'$, dividing gives the claim. Gronwall then needs $\int_0^T\norm{\nabla u}_{L^\infty} < \infty$; in three dimensions the case study gets this from $\nabla u \in L^1(0,T;W^{1,q})$ and $W^{1,q} \subset L^\infty$, valid only for $q > 3$.
  \end{solution}
\end{exercise}

\subsection*{Where to find a proof}

\begin{itemize}
\item Y. Cho, H. J. Choe, H. Kim, Unique solvability of the initial boundary value problems for compressible viscous fluids, \emph{J. Math. Pures Appl.} 83 (2004) 243--275, estimate (2.11).
\item R. J. DiPerna, P.-L. Lions, Ordinary differential equations, transport theory and Sobolev spaces, \emph{Invent. Math.} 98 (1989) 511--547 (non-smooth velocity fields; renormalization).
\end{itemize}

\neededfor


\chapter{What ``solution'' means}
% Strong, weak, and renormalized solutions

\section{Strong, weak, and renormalized solutions}\label{p:notions-of-solution}

\status{stub}
\prereqs{bochner-spaces, transport-equation}

\subsection*{The idea}

"Solution" means different things depending on how much regularity one can prove. Classical solutions have all derivatives continuous. Strong solutions have enough weak derivatives that the equation holds almost everywhere. Weak solutions satisfy the equation only after multiplying by test functions and integrating by parts. Each weaker notion is easier to construct and harder to make unique.

\subsection*{Statement}

\begin{defnbox}{Strong solution}
\index{strong solution}%
All derivatives appearing in the equation exist as functions (in some $L^p$), and the equation holds almost everywhere in $\Omega\times(0,T)$. The precise spaces are part of the definition; for example, the chemotaxis case study asks $\rho\in C([0,T];W^{1,q})$, $u\in C([0,T];D^2\cap D^1_0)\cap L^2(0,T;D^{2,q})$, $\sqrt\rho\,u_t\in L^\infty(0,T;L^2)$.
\end{defnbox}

\begin{defnbox}{Weak (distributional) solution}
\index{weak solution}%
The equation holds in $\mathcal{D}'$: multiplied by any smooth compactly supported test function and integrated by parts so that no derivative falls on the unknown beyond what its integrability allows. A \emph{finite-energy} weak solution additionally satisfies the energy inequality, which is part of the definition.
\end{defnbox}

\begin{defnbox}{Renormalized solution of the continuity equation}
\index{renormalized solution}%
$\rho$ solves $\rho_t+\dv(\rho u)=0$ in the renormalized sense if, for every $b\in C^1([0,\infty))$ with suitable growth (e.g. $b'$ bounded), $b(\rho)_t+\dv(b(\rho)u)+\big(b'(\rho)\rho-b(\rho)\big)\dv u=0$ in $\mathcal{D}'$.
\end{defnbox}

\begin{defnbox}{Compatibility condition}
\index{compatibility condition}%
A condition on the initial data saying that the equation, evaluated at $t=0$, is consistent with the regularity being asked for. In compressible flow with vacuum: if $\sqrt\rho\,u_t\in L^\infty(0,T;L^2)$ is required, then at $t=0$ the momentum equation forces $\Lame u_0-\nabla p(\rho_0)+(\text{forces}) = \sqrt{\rho_0}\,g$ with $g\in L^2$.
\end{defnbox}

\subsection*{Why it is plausible}

Renormalization is the chain rule: for smooth solutions, multiplying the equation by $b'(\rho)$ gives the renormalized form. For weak solutions this is not automatic, and asking for it is what allows nonlinear functions of $\rho$ (like the pressure) to be handled after weak limits. The compatibility condition is the same idea at $t=0$: the solution's regularity at the initial time cannot exceed what the data give, so the data must already have it.

\subsection*{Exercises}

\begin{exercise}
  Write the weak formulation of $u_t - \Delta u = f$ in $\Omega\times(0,T)$ with $u = 0$ on $\partial\Omega$ and $u(0) = u_0$. Show that a weak solution that is smooth up to the boundary is a classical solution.
  \begin{hint}
    Multiply by $\varphi \in C_c^\infty(\Omega\times[0,T))$, integrate by parts in $x$ and $t$; the initial value enters through $\int u_0\varphi(0)$. For a smooth weak solution, integrate back by parts: $\int(u_t - \Delta u - f)\varphi = 0$ for all $\varphi$ gives the equation, and then the remaining terms give the initial condition.
  \end{hint}
  \begin{solution}
    Multiply by $\varphi \in C_c^\infty(\Omega\times[0,T))$ and integrate. In time, $\int_0^T u_t\varphi = -\int_0^T u\varphi_t - u(0)\varphi(0)$; in space, $-\int\Delta u\,\varphi = \int\nabla u\cdot\nabla\varphi$. A weak solution is $u \in L^2(0,T;H^1_0)$ with

    \[ -\int_0^T\!\!\int_\Omega u\,\varphi_t + \int_0^T\!\!\int_\Omega\nabla u\cdot\nabla\varphi = \int_0^T\!\!\int_\Omega f\varphi + \int_\Omega u_0\,\varphi(\cdot,0)
             \quad\text{for all such }\varphi. \]

    If $u$ is smooth up to the boundary, integrate back by parts: $\int\!\!\int(u_t - \Delta u - f)\varphi + \int(u(\cdot,0) - u_0)\varphi(\cdot,0) = 0$. Taking $\varphi$ supported in $\Omega\times(0,T)$, the fundamental lemma gives $u_t - \Delta u = f$. Then $\int(u(\cdot,0) - u_0)\varphi(\cdot,0) = 0$ for all $\varphi$, so $u(0) = u_0$. Finally, $u(t) \in H^1_0$ and continuous up to $\partial\Omega$ means $u = 0$ there.
  \end{solution}
\end{exercise}

\begin{exercise}
  Renormalization. For smooth $\rho > 0$ and $u$ solving $\rho_t + \dv(\rho u) = 0$, verify the renormalized equation. Write it out for $b(\rho) = \rho\log\rho$; this is the form used in the Lions--Feireisl argument.
  \begin{hint}
    $b(\rho)_t + \dv(b(\rho)u) = b'(\rho)(\rho_t + u\cdot\nabla\rho) + b(\rho)\dv u = -\big(b'(\rho)\rho - b(\rho)\big)\dv u$. For $b = \rho\log\rho$: $b'(\rho)\rho - b(\rho) = \rho$, so $(\rho\log\rho)_t + \dv(\rho\log\rho\,u) + \rho\,\dv u = 0$.
  \end{hint}
  \begin{solution}
    Using $\rho_t = -\dv(\rho u) = -u\cdot\nabla\rho - \rho\,\dv u$:

    \[ b(\rho)_t + \dv(b(\rho)u) = b'(\rho)\big(\rho_t + u\cdot\nabla\rho\big) + b(\rho)\,\dv u = -\big(b'(\rho)\rho - b(\rho)\big)\dv u. \]

    For $b(\rho) = \rho\log\rho$, $b'(\rho) = 1 + \log\rho$ and $b'(\rho)\rho - b(\rho) = \rho$, so $(\rho\log\rho)_t + \dv(\rho\log\rho\,u) + \rho\,\dv u = 0$.
  \end{solution}
\end{exercise}

\begin{exercise}
  Weak solutions need not be unique. For Burgers' equation $u_t + (\frac12u^2)_x = 0$ with $u_0 = 0$, check that $u = -1$ on $-\frac t2 < x < 0$, $u = 1$ on $0 < x < \frac t2$, and $u = 0$ elsewhere is a weak solution. Show that its energy $\int\frac12u^2$ increases, so an energy inequality rules it out.
  \begin{hint}
    A piecewise constant function is a weak solution exactly when each jump moves with speed $s = \frac{[\frac12u^2]}{[u]}$ (Rankine--Hugoniot; derive it by testing across the jump). Check the three jumps: speeds $-\frac12$, $0$, $\frac12$. The energy is $\frac t2$.
  \end{hint}
  \begin{solution}
    \emph{Rankine--Hugoniot.} A function is a weak solution if $\int_0^\infty\!\int(u\varphi_t + \frac12u^2\varphi_x)\,dx\,dt + \int u_0\varphi(x,0)\,dx = 0$ for all test functions. On each region where $u$ is constant the equation holds classically; integrating by parts region by region, only the jump lines contribute. On a line $x = st$ the contribution is $\int(s[u] - [\frac12u^2])\varphi\,d\ell$ up to a positive factor, where $[\cdot]$ is right value minus left value. So $u$ is a weak solution iff every jump satisfies $s = [\frac12u^2]/[u]$.

    Check the three jumps: at $x = -\frac t2$, from $0$ to $-1$: $\frac{1/2 - 0}{-1 - 0} = -\frac12$. At $x = 0$, from $-1$ to $1$: $\frac{1/2 - 1/2}{2} = 0$. At $x = \frac t2$, from $1$ to $0$: $\frac{0 - 1/2}{0 - 1} = \frac12$. All match, and $u \to 0 = u_0$ as $t \to 0$.

    Its energy is $\int\frac12u^2\,dx = \frac12\cdot\frac t2 + \frac12\cdot\frac t2 = \frac t2$, which increases from $0$, while $u \equiv 0$ is another weak solution with the same data. Requiring that the energy not increase rules out the first.
  \end{solution}
\end{exercise}

\begin{exercise}
  Compatibility for the heat equation. Suppose $u_t = \Delta u$ with $u = 0$ on $\partial\Omega$, and $u_t$ is continuous up to $t = 0$ and to the boundary. Show that the data must satisfy $\Delta u_0 = 0$ on $\partial\Omega$.
  \begin{hint}
    $u = 0$ on $\partial\Omega$ for all $t$, so $u_t = 0$ there; at $t = 0$ the equation gives $u_t(0) = \Delta u_0$.
  \end{hint}
  \begin{solution}
    For $x \in \partial\Omega$, $u(x,t) = 0$ for all $t$, so $u_t(x,t) = 0$ for $t > 0$, and by continuity also at $t = 0$. By the equation and continuity at $t = 0$, $u_t(x,0) = \Delta u(x,0) = \Delta u_0(x)$. Hence $\Delta u_0 = 0$ on $\partial\Omega$. This is a condition on the data alone: smooth data that violate it cannot produce a solution with $u_t$ continuous up to $t = 0$.
  \end{solution}
\end{exercise}

\subsection*{Where to find a proof}

\begin{itemize}
\item E. Feireisl, A. Novotný, \emph{Singular Limits in Thermodynamics of Viscous Fluids}, Birkhäuser, 2009, §10.18 (renormalized solutions).
\item P.-L. Lions, \emph{Mathematical Topics in Fluid Mechanics, Vol. 2: Compressible Models}, Oxford, 1998.
\item Y. Cho, H. J. Choe, H. Kim, Unique solvability of the initial boundary value problems for compressible viscous fluids, \emph{J. Math. Pures Appl.} 83 (2004) 243--275 (strong solutions with vacuum and the compatibility condition).
\end{itemize}

\neededfor


\chapter{A priori estimates}
% Energy estimates and a priori bounds

\section{Energy estimates and a priori bounds}\label{p:energy-estimates}

\status{stub}
\prereqs{elementary-inequalities, gronwall}

\subsection*{The idea}

An a priori estimate is a bound that every smooth solution must satisfy, proved before knowing that solutions exist. The standard way to find one: multiply the equation by a well-chosen quantity, integrate, integrate by parts, and use the structure of the equation so that the dangerous terms cancel or have a sign. The choice of multiplier is the idea; the rest is the \pref{elementary-inequalities}{inequality toolkit}.

\subsection*{Statement}

\begin{thmbox}{The method (not a theorem)}
\index{energy estimate}%
\begin{enumerate}
\item Multiply each equation by a multiplier (often the unknown itself, its time derivative, or $\Delta$ of it) and integrate over $\Omega$. How to choose it is explained under \emph{Choosing the multiplier} below.
\item Integrate by parts, using the boundary conditions, to produce $\frac{d}{dt}$(energy) plus nonnegative dissipation.
\item Arrange the multipliers so that coupling terms cancel between equations.
\item Bound what remains by the energy and dissipation (Hölder, Sobolev, Young), absorb the dissipation parts, and close with \pref{gronwall}{Gronwall}.
\item For higher regularity, repeat after differentiating the equation in $t$ or $x$.
\end{enumerate}
\end{thmbox}

\begin{exbox}{Model example (heat equation)}
For $u_t-\Delta u=0$ with $u=0$ on $\partial\Omega$, multiplying by $u$ gives $\frac{d}{dt}\frac12\norm{u}_{L^2}^2+\norm{\nabla u}_{L^2}^2=0$, and multiplying by $u_t$ gives $\norm{u_t}_{L^2}^2+\frac{d}{dt}\frac12\norm{\nabla u}_{L^2}^2=0$.
\end{exbox}

\subsection*{Choosing the multiplier}

The multiplier is not guessed. First decide which quantity you want to bound; the multiplier is its derivative.

\begin{thmbox}{The multiplier rule}
\index{multiplier rule}\index{energy estimate!choosing the multiplier}%
To control $\int_\Omega \Phi(u)\,dx$, multiply the equation by $\Phi'(u)$. By the chain rule the time term becomes exactly the time derivative of that quantity:
\[
\int_\Omega \Phi'(u)\,u_t\,dx=\frac{d}{dt}\int_\Omega \Phi(u)\,dx.
\]
For the diffusion term, integration by parts gives
\[
-\int_\Omega \Delta u\,\Phi'(u)\,dx=\int_\Omega \Phi''(u)\,\abs{\nabla u}^2\,dx-\int_{\partial\Omega}\Phi'(u)\,\partial_\nu u\,dS .
\]
The boundary term vanishes under Neumann conditions ($\partial_\nu u=0$), or under Dirichlet conditions ($u=0$) when $\Phi'(0)=0$. If $\Phi$ is convex, $\Phi''\ge 0$, so the remaining term is nonnegative: it is dissipation, which can be kept or dropped.
\end{thmbox}

The standard choices for $u_t-\Delta u=f$:
\begin{itemize}
\item $\Phi(s)=\frac12 s^2$, multiplier $u$. Dissipation $\norm{\nabla u}_{L^2}^2$. This is the basic energy estimate.
\item $\Phi(s)=\frac1p\abs{s}^p$ ($p\ge 2$), multiplier $\abs{u}^{p-2}u$. Dissipation $(p-1)\int\abs{u}^{p-2}\abs{\nabla u}^2$. Controls $\norm{u}_{L^p}$.
\item $\Phi(s)=s\log s$ (for $u>0$), multiplier $\log u+1$. Dissipation $\int \abs{\nabla u}^2/u$. This is the entropy, the standard functional for chemotaxis (Keller--Segel) models.
\item $\Phi(s)=\abs{s}$, multiplier $\operatorname{sign}(u)$. Controls $\norm{u}_{L^1}$, natural for densities and mass. Since $\abs{s}$ is not differentiable at $0$, it has to be smoothed first (next example).
\end{itemize}

\begin{exbox}{The smoothed sign and the $L^1$ estimate}
Take $u_t-\Delta u=f$ with $u=0$ on $\partial\Omega$. Replace $\abs{s}$ by the smooth convex function
\[
\Phi_\delta(s)=\sqrt{s^2+\delta^2},\qquad
\Phi_\delta'(s)=\frac{s}{\sqrt{s^2+\delta^2}}\;\longrightarrow\;\operatorname{sign}(s),\qquad
\Phi_\delta''(s)=\frac{\delta^2}{(s^2+\delta^2)^{3/2}}\ge 0 .
\]
Since $\Phi_\delta'(0)=0$, the boundary term vanishes, and multiplying by $\Phi_\delta'(u)$ gives
\[
\frac{d}{dt}\int_\Omega \Phi_\delta(u)\,dx+\int_\Omega \Phi_\delta''(u)\abs{\nabla u}^2\,dx=\int_\Omega f\,\frac{u}{\sqrt{u^2+\delta^2}}\,dx\le \norm{f}_{L^1}.
\]
Drop the nonnegative term and integrate in time. Since $\abs{s}\le\Phi_\delta(s)\le\abs{s}+\delta$, letting $\delta\to 0$ gives
\[
\norm{u(t)}_{L^1}\le \norm{u_0}_{L^1}+\int_0^t\norm{f(s)}_{L^1}\,ds .
\]
Applied to the difference of two solutions of a linear problem, the same computation gives $L^1$ contraction, and hence uniqueness.
\end{exbox}

\paragraph{Multipliers with derivatives.} The same rule applies when the quantity you want involves derivatives of $u$:
\begin{itemize}
\item To control $\frac12\norm{\nabla u}_{L^2}^2$: its time derivative is $\int\nabla u\cdot\nabla u_t=\int(-\Delta u)\,u_t$, so the multiplier is $-\Delta u$. Equivalently, apply $\partial_i$ to the equation and multiply by $\partial_i u$: this is the basic estimate for $\nabla u$.
\item To control $\frac12\norm{\Delta u}_{L^2}^2$: the multiplier is $\Delta^2 u$.
\item The multiplier $u_t$ gives $\norm{u_t}_{L^2}^2+\frac{d}{dt}\frac12\norm{\nabla u}_{L^2}^2$. It controls the time derivative, which is what compactness (Aubin--Lions) needs.
\end{itemize}

\begin{rmkbox}{Reading off the function spaces}
Every estimate above has the form
\[
E(t)+\int_0^t D(s)\,ds\le E(0)\quad(\text{plus data terms}),
\]
with $E\ge 0$ the energy and $D\ge 0$ the dissipation. Each term of $E$ is bounded at every time, so it gives a space that is $L^\infty$ in time. Each term of $D$ is bounded only after integrating in time; since it appears squared, it gives a space that is $L^2$ in time. For the heat equation, the estimate at level $s$ (multiplier $u$, $-\Delta u$, $\Delta^2u$, \dots) gives
\[
u\in L^\infty(0,T;H^s)\cap L^2(0,T;H^{s+1}).
\]
The level is chosen by the nonlinearity, not by the estimate: you go up until the nonlinear terms can be bounded, typically until $s>d/2$, where $H^s\subset L^\infty$.
\end{rmkbox}

\subsection*{Why it is plausible}

Physical systems usually dissipate energy, and the mathematical energy is the physical one. Kinetic energy, internal (pressure) energy and potential energy together cannot grow; each piece then bounds a norm of the unknowns. When the energy contains a term with no sign, the first job is to show it is still bounded below, typically by interpolation against a conserved quantity.

\subsection*{Exercises}

\begin{exercise}
  Derive the two identities for the heat equation in the model example, with every integration by parts and boundary term. Conclude that $u_0 \in H^1_0$ gives $u \in L^\infty(0,T;H^1_0)$ and $u_t \in L^2(0,T;L^2)$.
  \begin{hint}
    $\int u_t u = \frac12\frac{d}{dt}\norm{u}^2$ and $-\int\Delta u\,u = \norm{\nabla u}^2$ (boundary term vanishes since $u = 0$). Similarly $-\int\Delta u\,u_t = \frac12\frac{d}{dt}\norm{\nabla u}^2$ because $u_t = 0$ on $\partial\Omega$. Integrate in time.
  \end{hint}
  \begin{solution}
    Multiplying by $u$: $\int u_tu = \frac12\frac{d}{dt}\norm{u}^2$ and $-\int\Delta u\,u = \int|\nabla u|^2 - \int_{\partial\Omega}u\,\partial_\nu u = \norm{\nabla u}^2$ since $u = 0$ on $\partial\Omega$. So $\frac12\frac{d}{dt}\norm{u}^2 + \norm{\nabla u}^2 = 0$.

    Multiplying by $u_t$: $-\int\Delta u\,u_t = \int\nabla u\cdot\nabla u_t - \int_{\partial\Omega}u_t\,\partial_\nu u = \frac12\frac{d}{dt}\norm{\nabla u}^2$, because $u_t = 0$ on $\partial\Omega$ ($u$ vanishes there for all $t$). So $\norm{u_t}^2 + \frac12\frac{d}{dt}\norm{\nabla u}^2 = 0$.

    Integrating: $\norm{u(t)}^2 \le \norm{u_0}^2$ and $\norm{\nabla u(t)}^2 + 2\int_0^t\norm{u_t}^2 = \norm{\nabla u_0}^2$. So $u \in L^\infty(0,T;H^1_0)$ and $u_t \in L^2(0,T;L^2)$ when $u_0 \in H^1_0$.
  \end{solution}
\end{exercise}

\begin{exercise}
  A nonlinear term with a sign. For $u_t - \Delta u = u - u^3$ with $\partial_\nu u = 0$, multiply by $u$ and show $\norm{u(t)}_{L^2}^2 \le \norm{u_0}_{L^2}^2 + \frac{|\Omega|}{2}t$.
  \begin{hint}
    $\frac12\frac{d}{dt}\norm{u}^2 + \norm{\nabla u}^2 = \int(u^2 - u^4)$, and $s^2 - s^4 \le \frac14$ for all real $s$.
  \end{hint}
  \begin{solution}
    Multiplying by $u$ and using $\partial_\nu u = 0$: $\frac12\frac{d}{dt}\norm{u}^2 + \norm{\nabla u}^2 = \int_\Omega(u^2 - u^4)$. Since $s^2 - s^4 = \frac14 - (s^2 - \frac12)^2 \le \frac14$, the right side is at most $\frac{|\Omega|}{4}$. Dropping $\norm{\nabla u}^2 \ge 0$, $\frac{d}{dt}\norm{u}^2 \le \frac{|\Omega|}{2}$, and integrating, $\norm{u(t)}^2 \le \norm{u_0}^2 + \frac{|\Omega|}{2}t$. The cubic term was not estimated at all: its sign did the work.
  \end{solution}
\end{exercise}

\begin{exercise}
  Choosing the multiplier. For the wave equation $u_{tt} = u_{xx}$ on $(0,1)$ with $u = 0$ at both ends, find a multiplier that shows $\int_0^1 (u_t^2 + u_x^2)\,dx$ is constant in time.
  \begin{hint}
    Multiply by $u_t$: $u_{tt}u_t = \frac12(u_t^2)_t$ and $\int u_{xx}u_t = -\int u_x u_{xt} = -\frac12\frac{d}{dt}\int u_x^2$, since $u_t = 0$ at the ends.
  \end{hint}
  \begin{solution}
    Multiply by $u_t$ and integrate over $(0,1)$. First, $\int u_{tt}u_t = \frac12\frac{d}{dt}\int u_t^2$. Second, $\int u_{xx}u_t = \big[u_xu_t\big]_0^1 - \int u_xu_{xt} = -\frac12\frac{d}{dt}\int u_x^2$, since $u_t = 0$ at $x = 0, 1$. So $\frac{d}{dt}\int_0^1(u_t^2 + u_x^2)\,dx = 0$.
  \end{solution}
\end{exercise}

\begin{exercise}
  The $L^p$ estimate. For $u_t - \Delta u = 0$ with $u = 0$ on $\partial\Omega$ and $p \ge 2$, show that $\norm{u(t)}_{L^p} \le \norm{u_0}_{L^p}$ for all $t \ge 0$.
  \begin{hint}
    Use the multiplier rule with $\Phi(s) = \frac1p\abs{s}^p$. Check that $\Phi'(0) = 0$ and $\Phi'' \ge 0$.
  \end{hint}
  \begin{solution}
    Multiply by $\Phi'(u) = \abs{u}^{p-2}u$, which vanishes on $\partial\Omega$. The time term is $\int \abs{u}^{p-2}u\,u_t = \frac1p\frac{d}{dt}\int\abs{u}^p$. The diffusion term is $-\int \Delta u\,\abs{u}^{p-2}u = \int \nabla u\cdot\nabla(\abs{u}^{p-2}u) = (p-1)\int\abs{u}^{p-2}\abs{\nabla u}^2 \ge 0$, using $\Phi''(s) = (p-1)\abs{s}^{p-2}$ (continuous since $p \ge 2$). So $\frac{d}{dt}\norm{u}_{L^p}^p \le 0$, and integrating, $\norm{u(t)}_{L^p}^p \le \norm{u_0}_{L^p}^p$.
  \end{solution}
\end{exercise}

\begin{exercise}
  A higher-order estimate. On the torus $\mathbb{T}^d$ (periodic boundary conditions, so integration by parts has no boundary terms), let $u_t - \Delta u = f$. Multiply by $\Delta^2 u$ and show
  \[
  \frac{d}{dt}\norm{\Delta u}_{L^2}^2 + \norm{\nabla\Delta u}_{L^2}^2 \le \norm{\nabla f}_{L^2}^2 .
  \]
  Which spaces does this put $\Delta u$ and $\nabla\Delta u$ in?
  \begin{hint}
    Integrate by parts twice in the time term and once in the other two. Finish with Cauchy--Schwarz and Young: $ab \le \frac12a^2 + \frac12b^2$.
  \end{hint}
  \begin{solution}
    Time term: $\int u_t\,\Delta^2u = \int \Delta u_t\,\Delta u = \frac12\frac{d}{dt}\norm{\Delta u}^2$. Diffusion term: $-\int\Delta u\,\Delta^2u = \int\nabla\Delta u\cdot\nabla\Delta u = \norm{\nabla\Delta u}^2$. Forcing: $\int f\,\Delta^2u = -\int\nabla f\cdot\nabla\Delta u \le \norm{\nabla f}\norm{\nabla\Delta u} \le \frac12\norm{\nabla f}^2 + \frac12\norm{\nabla\Delta u}^2$. Absorbing the last term into the left side and multiplying by $2$ gives the inequality. Integrating in time, $\Delta u \in L^\infty(0,T;L^2)$ and $\nabla\Delta u \in L^2(0,T;L^2)$ whenever $\Delta u_0 \in L^2$ and $\nabla f \in L^2(0,T;L^2)$: the level-2 spaces $L^\infty(0,T;H^2)\cap L^2(0,T;H^3)$, up to lower-order terms.
  \end{solution}
\end{exercise}

\subsection*{Where to find a proof}

\begin{itemize}
\item L. C. Evans, \emph{Partial Differential Equations}, 2nd ed., AMS, 2010, §2.3.4 and §7.1.2 (energy methods for the heat equation and for linear parabolic equations).
\item J. C. Robinson, J. L. Rodrigo, W. Sadowski, \emph{The Three-Dimensional Navier--Stokes Equations: Classical Theory}, Cambridge University Press, 2016 (higher-order energy estimates written out line by line).
\item J. L. Vázquez, \emph{The Porous Medium Equation: Mathematical Theory}, Oxford University Press, 2007 ($L^1$ estimates and $L^1$ contraction with smoothed sign functions).
\end{itemize}

\neededfor


\chapter{Existence by approximation}
% The Faedo–Galerkin method

\section{The Faedo--Galerkin method}\label{p:galerkin}

\status{stub}
\prereqs{energy-estimates, weak-convergence, bochner-spaces}

\subsection*{The idea}

To solve a PDE, solve a finite-dimensional version of it (an ODE system), prove bounds that do not depend on the dimension, and pass to the limit. The a priori estimates are what make this work: the approximations are built so that the same estimates hold for them.

\subsection*{Statement}

\begin{thmbox}{The scheme}
\begin{enumerate}
\item Choose a basis $\{w_k\}$ of the solution space, usually eigenfunctions of the elliptic part of the equation (so the boundary conditions are built in).
\item Look for $u_m(t)=\sum_{k=1}^m c_k(t)w_k$ satisfying the weak form tested against $w_1,\dots,w_m$ only. This is an ODE system for $c_1,\dots,c_m$; solve it locally by ODE theory.
\item Prove the a priori estimates for $u_m$ with constants independent of $m$; they extend the ODE solution to a uniform time interval.
\item Extract a weakly convergent subsequence and show the limit satisfies the equation, using compactness for the nonlinear terms.
\end{enumerate}

For linear parabolic equations this gives existence and uniqueness of weak solutions with the energy regularity.
\end{thmbox}

\subsection*{Why it is plausible}

Testing the approximate equation with $u_m$ itself is allowed, because $u_m$ is a combination of the test functions. So the energy identity holds exactly at the finite level, with the same constants. In nonlinear problems the ODE system has a coefficient matrix (for compressible flow, $(\rho w_i,w_k)$, a weighted mass matrix) that must be invertible; this is one reason to start with density bounded away from zero.

\subsection*{Exercises}

\begin{exercise}
  Galerkin for the heat equation on $(0,\pi)$. With $w_k = \sin(kx)$ and $u_m = \sum_{k\le m}c_k(t)w_k$, write the Galerkin equations for $u_t = u_{xx} + f$, $u = 0$ at the ends. Solve them, and observe that $u_m$ is the truncated Fourier series of the exact solution.
  \begin{hint}
    Testing against $w_k$ decouples: $c_k' = -k^2c_k + f_k$ with $f_k = \frac2\pi\int_0^\pi f\,w_k$. The basis consists of eigenfunctions, so nothing couples different modes.
  \end{hint}
  \begin{solution}
    The weak form is $(u_t, v) + (u_x, v_x) = (f, v)$ for $v \in H^1_0(0,\pi)$. On $(0,\pi)$, $(w_j, w_k) = \frac\pi2\delta_{jk}$ and $(w_j', w_k') = \frac\pi2k^2\delta_{jk}$. With $u_m = \sum_{k\le m}c_kw_k$ and $v = w_j$, $j \le m$:

    \[ \frac\pi2c_j' + \frac\pi2j^2c_j = (f, w_j), \quad\text{i.e.}\quad c_j' = -j^2c_j + f_j,\ \ f_j = \frac2\pi(f,w_j), \quad c_j(0) = \frac2\pi(u_0,w_j). \]

    So $c_j(t) = e^{-j^2t}c_j(0) + \int_0^te^{-j^2(t-s)}f_j(s)\,ds$. Testing the exact equation with $w_j$ shows that the $j$-th sine coefficient of the exact solution solves the same ODE with the same initial value. Hence $u_m$ is the exact solution's Fourier series cut at $m$: an eigenfunction basis decouples the modes completely.
  \end{solution}
\end{exercise}

\begin{exercise}
  Show that the Galerkin equations $(\partial_t u_m, w_k) + (\nabla u_m, \nabla w_k) = (f, w_k)$, $k \le m$, imply the energy identity $\frac12\frac{d}{dt}\norm{u_m}^2 + \norm{\nabla u_m}^2 = (f, u_m)$, with constants independent of $m$.
  \begin{hint}
    Multiply the $k$-th equation by $c_k(t)$ and sum over $k$. This is "testing with $u_m$", allowed because $u_m$ is a combination of the test functions.
  \end{hint}
  \begin{solution}
    Multiply the $k$-th equation by $c_k(t)$ and sum over $k \le m$. By linearity, $(\partial_tu_m, \sum_kc_kw_k) + (\nabla u_m, \nabla\sum_kc_kw_k) = (f, \sum_kc_kw_k)$, i.e. $(\partial_tu_m, u_m) + \norm{\nabla u_m}^2 = (f, u_m)$, and $(\partial_tu_m, u_m) = \frac12\frac{d}{dt}\norm{u_m}^2$. No constant depends on $m$. Every estimate derived from this identity, for instance with Young and Gronwall, is therefore uniform in $m$.
  \end{solution}
\end{exercise}

\begin{exercise}
  Passing to the limit. If $u_m \weakto u$ in $L^2(0,T;H^1_0)$, show $\int_0^T(\nabla u_m,\nabla w)\psi\,dt \to \int_0^T(\nabla u,\nabla w)\psi\,dt$ for fixed $w \in H^1_0$ and $\psi \in C_c^\infty(0,T)$. Explain why $\int_0^T\!\int u_m^2 w\,\psi$ need not converge to the same expression with $u$.
  \begin{hint}
    The map $v \mapsto \int_0^T(\nabla v,\nabla w)\psi$ is a bounded linear functional on $L^2(0,T;H^1_0)$. Squares are not linear: see the $\sin(nx)$ example in \pref{weak-convergence}{weak and weak* convergence}.
  \end{hint}
  \begin{solution}
    $\ell(v) = \int_0^T(\nabla v, \nabla w)\psi\,dt$ is linear and $|\ell(v)| \le \norm{\nabla w}_{L^2}\norm{\psi}_{L^2(0,T)}\norm{v}_{L^2(0,T;H^1_0)}$, so $\ell$ is a bounded linear functional on $L^2(0,T;H^1_0)$ and $\ell(u_m) \to \ell(u)$ by definition of weak convergence.

    The map $v \mapsto \int_0^T\!\!\int v^2w\psi$ is not linear, and weak convergence does not pass through it. Example: fix $v \in H^1_0$, $v \ne 0$, and $u_m(x,t) = \sin(mt)v(x)$. Then $u_m \weakto 0$ in $L^2(0,T;H^1_0)$, but $\int_0^T\!\!\int u_m^2w\psi = \int_0^T\sin^2(mt)\psi(t)\,dt\int v^2w \to \frac12\int_0^T\psi\int v^2w$, which is not $0$ in general.
  \end{solution}
\end{exercise}

\subsection*{Where to find a proof}

\begin{itemize}
\item L. C. Evans, \emph{Partial Differential Equations}, 2nd ed., AMS, 2010, §7.1.2 (Theorems 1--3, linear parabolic equations).
\item M. Padula, Existence of global solutions for 2-dimensional viscous compressible flows, \emph{J. Funct. Anal.} 69 (1986) 1--20.
\item H. J. Choe, H. Kim, Strong solutions of the Navier--Stokes equations for isentropic compressible fluids, \emph{J. Differential Equations} 190 (2003) 504--523.
\end{itemize}

\neededfor


\chapter{Uniqueness}
% Uniqueness by difference estimates

\section{Uniqueness by difference estimates}\label{p:uniqueness-energy}

\status{stub}
\prereqs{energy-estimates, gronwall}

\subsection*{The idea}

To show two solutions with the same data coincide, subtract them and estimate the difference. The trick is to measure the difference in a weaker norm than the solutions themselves live in, so that the regularity of the two solutions can be used freely as coefficients.

\subsection*{Statement}

\begin{thmbox}{The scheme}
\begin{enumerate}
\item Write the equations for the difference $(\bar\rho,\bar u,\bar\phi)$ of two solutions. Nonlinear terms become linear in the difference, with coefficients built from the two solutions.
\item Run an energy estimate for the difference at a low level (typically $L^2$), producing $\frac{d}{dt}Y + \text{dissipation}\le K(t)\,Y$.
\item Check $K\in L^1(0,T)$ using the regularity of the solutions, and apply \pref{gronwall}{Gronwall} with $Y(0)=0$.
\end{enumerate}

A variant: \emph{weak--strong uniqueness} compares a weak solution with a strong one with the same data, using the energy inequality of the weak solution in place of the equation, and shows they coincide.
\end{thmbox}

\subsection*{Why it is plausible}

Uniqueness asks less than existence: one does not need new bounds, only that the bounds already available are strong enough to control the coefficient $K(t)$. That is why uniqueness theorems are usually stated in the same class where existence of strong solutions was proved, and why uniqueness of weak solutions is often open.

\subsection*{Exercises}

\begin{exercise}
  Let $u_1, u_2 \in L^\infty(\Omega\times(0,T))$ be solutions of $u_t - \Delta u = u^2$ with the same data and $u = 0$ on $\partial\Omega$. Show $u_1 = u_2$ by an $L^2$ estimate for $w = u_1 - u_2$.
  \begin{hint}
    $w_t - \Delta w = (u_1 + u_2)w$. Testing with $w$: $\frac12\frac{d}{dt}\norm{w}^2 + \norm{\nabla w}^2 \le \norm{u_1 + u_2}_{L^\infty}\norm{w}^2$. Gronwall with $w(0) = 0$; here $K(t) = 2\norm{u_1 + u_2}_{L^\infty}$.
  \end{hint}
  \begin{solution}
    $w = u_1 - u_2$ satisfies $w_t - \Delta w = u_1^2 - u_2^2 = (u_1 + u_2)w$, $w = 0$ on $\partial\Omega$, $w(0) = 0$. Testing with $w$:

    \[ \frac12\frac{d}{dt}\norm{w}^2 + \norm{\nabla w}^2 = \int(u_1 + u_2)w^2 \le \norm{u_1 + u_2}_{L^\infty}\norm{w}^2. \]

    So $Y = \norm{w}^2$ satisfies $Y' \le KY$ with $K = 2\norm{u_1 + u_2}_{L^\infty} \in L^1(0,T)$, and $Y(0) = 0$. By Gronwall, $Y \equiv 0$. The boundedness of the two solutions was used only as a coefficient.
  \end{solution}
\end{exercise}

\begin{exercise}
  Viscous Burgers on $(0,1)$: $u_t + uu_x = u_{xx}$, $u = 0$ at both ends. For two solutions, find the coefficient $K(t)$ in the $L^2$ estimate for the difference, and say which norms of $u_1, u_2$ must be integrable in time.
  \begin{hint}
    For $w = u_1 - u_2$: $w_t + u_1w_x + u_{2,x}w = w_{xx}$. Bound $|\int u_1ww_x| \le \frac12\norm{w_x}^2 + \frac12\norm{u_1}_{L^\infty}^2\norm{w}^2$ and $|\int u_{2,x}w^2| \le \norm{u_{2,x}}_{L^\infty}\norm{w}^2$. So $K = \norm{u_1}_{L^\infty}^2 + 2\norm{u_{2,x}}_{L^\infty}$ must be in $L^1(0,T)$.
  \end{hint}
  \begin{solution}
    Since $u_1u_{1,x} - u_2u_{2,x} = u_1w_x + wu_{2,x}$, the difference satisfies $w_t + u_1w_x + u_{2,x}w = w_{xx}$. Testing with $w$:

    \[ \frac12\frac{d}{dt}\norm{w}^2 + \norm{w_x}^2 = -\int u_1ww_x - \int u_{2,x}w^2. \]

    Bound $|\int u_1ww_x| \le \norm{u_1}_{L^\infty}\norm{w}\norm{w_x} \le \frac12\norm{w_x}^2 + \frac12\norm{u_1}_{L^\infty}^2\norm{w}^2$ and $|\int u_{2,x}w^2| \le \norm{u_{2,x}}_{L^\infty}\norm{w}^2$. Absorbing, $\frac{d}{dt}\norm{w}^2 \le \big(\norm{u_1}_{L^\infty}^2 + 2\norm{u_{2,x}}_{L^\infty}\big)\norm{w}^2$. So uniqueness holds whenever $u_1 \in L^2(0,T;L^\infty)$ and $u_{2,x} \in L^1(0,T;L^\infty)$.
  \end{solution}
\end{exercise}

\subsection*{Where to find a proof}

\begin{itemize}
\item L. C. Evans, \emph{Partial Differential Equations}, 2nd ed., AMS, 2010, §7.1.2 (Theorem 4, uniqueness for linear parabolic equations).
\item X. Huo, A. Jüngel, Global existence and weak-strong uniqueness for chemotaxis compressible Navier--Stokes equations modeling vascular network formation, \emph{J. Math. Fluid Mech.} 26 (2024) 11.
\end{itemize}

\neededfor


\chapter{Global weak solutions}
% Compactness: Rellich–Kondrachov and Aubin–Lions

\section{Compactness: Rellich--Kondrachov and Aubin--Lions}\label{p:compactness}

\status{stub}
\prereqs{sobolev-embedding, bochner-spaces, weak-convergence}

\subsection*{The idea}

Weak convergence cannot pass limits through nonlinear terms, so one needs strong convergence of at least one factor. Compactness theorems give it: a bound on derivatives makes a bounded set precompact in a weaker norm. In space this is Rellich--Kondrachov; in space--time it is Aubin--Lions, which needs a bound on the time derivative as well.

\subsection*{Statement}

\begin{thmbox}{Rellich--Kondrachov}
\index{Rellich--Kondrachov theorem}\index{compactness!Rellich--Kondrachov}%
Let $\Omega\subset\R^n$ be bounded with Lipschitz boundary and $1\le p< n$. Then $W^{1,p}(\Omega)$ embeds compactly in $L^q(\Omega)$ for every $1\le q< p^*$. For $p\ge n$ the embedding into $L^q$ is compact for every $q< \infty$.
\end{thmbox}

\begin{thmbox}{Aubin--Lions (Simon's form)}
\index{Aubin--Lions lemma}\index{compactness!Aubin--Lions}%
Let $X_0\subset X\subset X_1$ be Banach spaces with $X_0\subset X$ compact and $X\subset X_1$ continuous. Let $F$ be bounded in $L^p(0,T;X_0)$ with $\{\partial_t f: f\in F\}$ bounded in $L^r(0,T;X_1)$.

\begin{itemize}
\item If $1\le p< \infty$ and $r\ge1$, then $F$ is relatively compact in $L^p(0,T;X)$.
\item If $p=\infty$ and $r> 1$, then $F$ is relatively compact in $C([0,T];X)$.
\end{itemize}
\end{thmbox}

\subsection*{Why it is plausible}

A bound on $\nabla u$ prevents fast oscillation in space, and fast oscillation is exactly what weak convergence loses (the $\sin(nx)$ example on the \pref{weak-convergence}{weak convergence} page has $\norm{\nabla u_n}\to\infty$). A bound on $u_t$, even in a very weak space, prevents oscillation in time. With both, no oscillation is left, and weak convergence improves to strong. The art is in finding which quantity has a time derivative bound: for compressible flow it is the momentum $\rho u$, not $u$.

\subsection*{Exercises}

\begin{exercise}
  Rellich in one dimension. Let $(u_n)$ be bounded in $H^1_0(0,1)$, $u_n = \sum_k a_k^{(n)}\sqrt2\sin(k\pi x)$. Show $\sum_{k> N}|a_k^{(n)}|^2 \le \frac{M}{N^2\pi^2}$ uniformly in $n$, and use this to extract a subsequence converging in $L^2$.
  \begin{hint}
    $\norm{u_n'}_{L^2}^2 = \sum_k k^2\pi^2|a_k^{(n)}|^2 \le M$. Choose a subsequence along which each coefficient converges (diagonal argument); the tails are uniformly small, so the subsequence is Cauchy in $L^2$.
  \end{hint}
  \begin{solution}
    The functions $\sqrt2\sin(k\pi x)$ and $\sqrt2\cos(k\pi x)$, $k \ge 1$, are orthonormal in $L^2(0,1)$. For $u \in H^1_0$, integration by parts (using $u(0) = u(1) = 0$) gives $(u', \sqrt2\cos(k\pi x)) = k\pi a_k$, so Bessel's inequality yields $\sum_kk^2\pi^2a_k^2 \le \norm{u'}^2$. If $\norm{u_n'}^2 \le M$ for all $n$, then

    \[ \sum_{k> N}|a_k^{(n)}|^2 \le \frac{1}{N^2\pi^2}\sum_kk^2\pi^2|a_k^{(n)}|^2 \le \frac{M}{N^2\pi^2}. \]

    Each sequence $(a_k^{(n)})_n$ is bounded, so a diagonal argument gives a subsequence along which $a_k^{(n)}$ converges for every $k$. Along it, $\norm{u_n - u_m}_{L^2}^2 \le \sum_{k\le N}|a_k^{(n)} - a_k^{(m)}|^2 + \frac{4M}{N^2\pi^2}$. Choose $N$ to make the last term small, then $n, m$ large: the subsequence is Cauchy in $L^2$, hence convergent.
  \end{solution}
\end{exercise}

\begin{exercise}
  Why a time-derivative bound is needed. Let $w \in H^1_0(\Omega)$, $w \ne 0$, and $u_n(x,t) = \sin(nt)w(x)$. Show $(u_n)$ is bounded in $L^2(0,T;H^1_0)$ but has no subsequence converging in $L^2(0,T;L^2)$. What happens to $\partial_t u_n$?
  \begin{hint}
    In time this is the $\sin(nt)$ example: weak limit $0$, norm bounded below. $\partial_t u_n = n\cos(nt)w$ is unbounded in every $L^r(0,T;X_1)$.
  \end{hint}
  \begin{solution}
    $\norm{u_n}_{L^2(0,T;H^1_0)}^2 = \norm{w}_{H^1}^2\int_0^T\sin^2(nt)\,dt \le T\norm{w}_{H^1}^2$: bounded. For $g \in L^2(0,T;L^2)$, $\int_0^T\!\!\int u_ng = \int_0^T\sin(nt)(w, g(t))\,dt \to 0$ (the $\sin(nx)$ argument in time), so any strong limit of a subsequence would be $0$. But $\norm{u_n}_{L^2(0,T;L^2)}^2 = \norm{w}_{L^2}^2\int_0^T\sin^2(nt)\,dt \to \frac T2\norm{w}_{L^2}^2 \ne 0$. Meanwhile $\partial_tu_n = n\cos(nt)w$ has norm of order $n$ in every $L^r(0,T;X_1)$: the hypothesis of Aubin--Lions fails exactly here.
  \end{solution}
\end{exercise}

\begin{exercise}
  Compactness fails at the critical exponent. For $u \in C_c^\infty(B_1)$, $p < n$, let $u_\lambda(x) = \lambda^{(n-p)/p}u(\lambda x)$. Show $\norm{\nabla u_\lambda}_{L^p}$ and $\norm{u_\lambda}_{L^{p^*}}$ do not depend on $\lambda$, and conclude that $W^{1,p}_0(B_1) \subset L^{p^*}(B_1)$ is not compact.
  \begin{hint}
    Change variables. As $\lambda \to \infty$ the supports shrink to the origin, so any limit in $L^{p^*}$ would be $0$, but the norms stay equal to $\norm{u}_{L^{p^*}} \ne 0$.
  \end{hint}
  \begin{solution}
    $\nabla u_\lambda(x) = \lambda^{\frac{n-p}{p}+1}(\nabla u)(\lambda x) = \lambda^{n/p}(\nabla u)(\lambda x)$, so $\norm{\nabla u_\lambda}_{L^p}^p = \lambda^n\lambda^{-n}\norm{\nabla u}_{L^p}^p$. With $p^* = \frac{np}{n-p}$, $\norm{u_\lambda}_{L^{p^*}}^{p^*} = \lambda^{\frac{n-p}{p}p^*}\lambda^{-n}\norm{u}_{L^{p^*}}^{p^*} = \lambda^n\lambda^{-n}\norm{u}_{L^{p^*}}^{p^*}$. Both are independent of $\lambda$.

    For $\lambda \ge 1$, $\operatorname{supp}u_\lambda \subset B_{1/\lambda}$, so $(u_\lambda)$ is bounded in $W^{1,p}_0(B_1)$. If the embedding were compact, some $u_{\lambda_k}$ with $\lambda_k \to \infty$ would converge in $L^{p^*}$, and a further subsequence a.e., to some $g$. For each $x \ne 0$, $u_{\lambda_k}(x) = 0$ once $\lambda_k > 1/|x|$, so $g = 0$; but $\norm{u_{\lambda_k}}_{L^{p^*}} = \norm{u}_{L^{p^*}} \ne 0$. Contradiction.
  \end{solution}
\end{exercise}

\subsection*{Where to find a proof}

\begin{itemize}
\item L. C. Evans, \emph{Partial Differential Equations}, 2nd ed., AMS, 2010, §5.7 (Rellich--Kondrachov).
\item J. Simon, Compact sets in the space $L^p(0,T;B)$, \emph{Ann. Mat. Pura Appl.} 146 (1987) 65--96, Corollary 4.
\end{itemize}

\neededfor

% Regularized approximation schemes

\section{Regularized approximation schemes}\label{p:regularized-approximation}

\status{stub}
\prereqs{energy-estimates, galerkin}

\subsection*{The idea}

When a system is too degenerate to solve directly, add small terms that make it solvable, prove estimates uniform in the small parameters, and remove the terms one at a time. Each added term must keep the energy structure intact, otherwise the uniform estimates are lost. Designing the regularization is often where the new idea of a paper lives.

\subsection*{Statement}

\begin{thmbox}{The standard scheme for compressible Navier--Stokes (Feireisl--Novotný--Petzeltová)}
\[
      \begin{gathered}
      \rho_t+\dv(\rho u)=\varepsilon\Delta\rho,\\
      (\rho u)_t+\dv(\rho u\otimes u)+\nabla p(\rho)+\delta\nabla\rho^\beta+\varepsilon\nabla\rho\cdot\nabla u
       =\mu\Delta u+(\lambda+\mu)\nabla\dv u,
      \end{gathered}
    \]

\begin{itemize}
\item $\varepsilon\Delta\rho$ (artificial viscosity) makes the continuity equation parabolic: $\rho$ becomes smooth and positive.
\item $\varepsilon\nabla\rho\cdot\nabla u$ cancels the energy produced by $\varepsilon\Delta\rho$, so the energy law survives.
\item $\delta\rho^\beta$ with $\beta$ large (artificial pressure) gives extra integrability of $\rho$ for the intermediate limits.
\end{itemize}

Limits are taken in order: Galerkin index $\to\infty$, then $\varepsilon\to0$, then $\delta\to0$.
\end{thmbox}

\subsection*{Why it is plausible}

Each regularization fixes one defect and costs one limit. The order matters: $\varepsilon\to0$ is done while $\delta\rho^\beta$ still provides integrability, and $\delta\to0$ last, when only the genuine pressure is left. In a coupled system, a regularization of one equation can break the energy balance with another, and then a new term must be invented to restore it.

\subsection*{Exercises}

\begin{exercise}
  For $\rho_t + \dv(\rho u) = \varepsilon\Delta\rho$ with $u = 0$ and $\partial_\nu\rho = 0$ on $\partial\Omega$, show that the mass is conserved. With $P(\rho) = \frac{\rho^\gamma}{\gamma-1}$ and $p(\rho) = \rho^\gamma$, show $\frac{d}{dt}\int P(\rho) + \varepsilon\int P''(\rho)|\nabla\rho|^2 = -\int p(\rho)\,\dv u$. The artificial viscosity only adds dissipation.
  \begin{hint}
    Multiply by $P'(\rho)$ and integrate. The viscosity term gives $-\varepsilon\int P''(\rho)|\nabla\rho|^2$. For the transport term, $\rho P''(\rho) = p'(\rho)$, so $\int P'(\rho)\dv(\rho u) = -\int\rho P''(\rho)\nabla\rho\cdot u = -\int\nabla p\cdot u = \int p\,\dv u$.
  \end{hint}
  \begin{solution}
    Mass: $\frac{d}{dt}\int\rho = \varepsilon\int_{\partial\Omega}\partial_\nu\rho - \int_{\partial\Omega}\rho\,u\cdot\nu = 0$.

    Multiply the equation by $P'(\rho)$ and integrate. With $\partial_\nu\rho = 0$, $\varepsilon\int P'(\rho)\Delta\rho = -\varepsilon\int P''(\rho)|\nabla\rho|^2$. Since $\rho P''(\rho) = \gamma\rho^{\gamma-1} = p'(\rho)$ and $u = 0$ on $\partial\Omega$,

    \[ \int P'(\rho)\dv(\rho u) = -\int P''(\rho)\nabla\rho\cdot\rho u = -\int\nabla p(\rho)\cdot u = \int p(\rho)\dv u. \]

    So $\frac{d}{dt}\int P(\rho) + \varepsilon\int P''(\rho)|\nabla\rho|^2 = -\int p(\rho)\dv u$. The right side is the pressure work that cancels against the momentum equation; the new term $\varepsilon\int P''|\nabla\rho|^2 \ge 0$ is extra dissipation.
  \end{solution}
\end{exercise}

\begin{exercise}
  The balancing term. Show $\frac{\varepsilon}{2}\int\Delta\rho\,|u|^2 = -\varepsilon\int(\nabla\rho\cdot\nabla)u\cdot u$ when $u = 0$ on $\partial\Omega$. This is why adding $\varepsilon\nabla\rho\cdot\nabla u$ to the momentum equation cancels the extra kinetic energy created by $\varepsilon\Delta\rho$.
  \begin{hint}
    Integrate by parts once: $\int\Delta\rho\,|u|^2 = -\int\nabla\rho\cdot\nabla|u|^2$, and $\partial_j|u|^2 = 2u_i\partial_ju_i$.
  \end{hint}
  \begin{solution}
    Integrating by parts once, $\int\Delta\rho\,|u|^2 = -\int\nabla\rho\cdot\nabla|u|^2 = -2\int\partial_j\rho\,u_i\partial_ju_i = -2\int(\nabla\rho\cdot\nabla)u\cdot u$, with no boundary term since $u = 0$ on $\partial\Omega$. Multiply by $\frac\varepsilon2$.

    Why this matters: testing $(\rho u)_t + \dv(\rho u\otimes u)$ with $u$ gives $\frac{d}{dt}\int\frac12\rho|u|^2 + \frac12\int\big(\rho_t + \dv(\rho u)\big)|u|^2$. In the original system the bracket vanishes; in the regularized one it equals $\varepsilon\Delta\rho$, leaving $\frac\varepsilon2\int\Delta\rho|u|^2$, which has no sign. Testing the balancing term $\varepsilon\nabla\rho\cdot\nabla u$ with $u$ gives $\varepsilon\int(\nabla\rho\cdot\nabla)u\cdot u = -\frac\varepsilon2\int\Delta\rho|u|^2$ by the identity. The two cancel exactly, and the kinetic energy balance is the same as without regularization.
  \end{solution}
\end{exercise}

\subsection*{Where to find a proof}

\begin{itemize}
\item E. Feireisl, A. Novotný, H. Petzeltová, On the existence of globally defined weak solutions to the Navier--Stokes equations, \emph{J. Math. Fluid Mech.} 3 (2001) 358--392.
\item A. Novotný, I. Straškraba, \emph{Introduction to the Mathematical Theory of Compressible Flow}, Oxford, 2004, Chapter 7.
\end{itemize}

\neededfor

% Lions–Feireisl theory for compressible flow
% Converted from notes/pde/lions-feireisl.html. Edit freely: this file is now the master copy.

\section{Lions--Feireisl theory for compressible flow}\label{p:lions-feireisl}

\prereqs{Compactness: Rellich--Kondrachov and Aubin--Lions~(\S\ref{p:compactness}), Strong, weak, and renormalized solutions~(\S\ref{p:notions-of-solution}), Regularized approximation schemes~(\S\ref{p:regularized-approximation})}

\subsection*{The idea}

The central difficulty for weak solutions of compressible Navier--Stokes: the density converges only weakly, so the weak limit of the pressure $p(\rho_n)$ need not be $p(\rho)$. Lions found a quantity that does behave well under weak limits, the effective viscous flux, and used it to prove strong convergence of the density. Feireisl extended the range of exponents to $\gamma> 3/2$.

\subsection*{Statement}

\begin{thmbox}{Theorem (Feireisl--Novotný--Petzeltová, 2001)}
Let $\Omega\subset\R^3$ be bounded with smooth boundary and $p(\rho)=a\rho^\gamma$ with $\gamma> 3/2$. For initial data of finite energy, the barotropic compressible Navier--Stokes equations with no-slip boundary condition have a global finite-energy weak solution, and the density satisfies the continuity equation in the renormalized sense. (Lions had proved this for $\gamma\ge9/5$.)
\end{thmbox}

\begin{rmkbox}{The ingredients}
\begin{itemize}
\item \emph{Effective viscous flux.} $F = (\lambda+2\mu)\dv u - p(\rho)$. Although neither term converges strongly, products like $T(\rho_n)F_n$ converge weakly to $\overline{T(\rho)}\,F$: the oscillations of $p(\rho)$ and of $\dv u$ cancel.
\item \emph{Renormalization} (see notions of solution~(\S\ref{p:notions-of-solution})), applied to truncations $T_k(\rho)$ of the density when $\rho$ is not square integrable.
\item \emph{Oscillation defect measure.} A bound on how much the truncations can oscillate; finite exactly when $\gamma> 3/2$.
\end{itemize}
\end{rmkbox}

\subsection*{Why it is plausible}

Strong convergence of $\rho_n$ is equivalent to $\overline{\rho\log\rho}=\rho\log\rho$ (weak limit equals function of the limit), and by convexity one always has $\ge$. The effective viscous flux identity, combined with the renormalized equations, shows that the defect $\overline{\rho\log\rho}-\rho\log\rho$ satisfies a transport inequality with zero initial data, so it stays zero. Below $\gamma=3/2$ the integrability of $\rho$ is too low for the products in this argument to make sense, even after truncation.

\subsection*{Exercises}

\begin{exercise}
  Let $\rho_n \ge 0$, $\rho_n \weakto \rho$ and $\rho_n\log\rho_n \weakto \overline{\rho\log\rho}$ in $L^1(\Omega)$, with $\rho > 0$. Show $\overline{\rho\log\rho} \ge \rho\log\rho$ a.e.
  \begin{hint}
    Convexity: $s\log s \ge \rho\log\rho + (1 + \log\rho)(s - \rho)$ for all $s \ge 0$. Put $s = \rho_n$, multiply by a nonnegative test function, and pass to the weak limit (the right-hand side is linear in $\rho_n$; truncate $\log\rho$ if it is unbounded).
  \end{hint}
  \begin{solution}
    $g(s) = s\log s$ is convex on $[0,\infty)$ ($g'' = 1/s > 0$), so it lies above its tangent at $\rho > 0$: $\rho_n\log\rho_n \ge \rho\log\rho + (1 + \log\rho)(\rho_n - \rho)$ pointwise. Fix $k$ and let $\psi \ge 0$ be bounded and supported in $\{1/k \le \rho \le k\}$, so that $(1 + \log\rho)\psi \in L^\infty$. Multiplying by $\psi$ and integrating,

    \[ \int\rho_n\log\rho_n\,\psi \ge \int\rho\log\rho\,\psi + \int(1 + \log\rho)\psi\,(\rho_n - \rho). \]

    Weak convergence in $L^1$ against $L^\infty$ functions gives, in the limit, $\int\overline{\rho\log\rho}\,\psi \ge \int\rho\log\rho\,\psi$. As this holds for all such $\psi$, $\overline{\rho\log\rho} \ge \rho\log\rho$ a.e. on $\{1/k \le \rho \le k\}$; let $k \to \infty$.
  \end{solution}
\end{exercise}

\begin{exercise}
  Strong convergence from "limit of the square equals square of the limit". On a bounded $\Omega$, show that $\rho_n \weakto \rho$ and $\rho_n^2 \weakto \rho^2$ in $L^1$ (with $\rho_n$ bounded in $L^2$) imply $\rho_n \to \rho$ in $L^2$.
  \begin{hint}
    $\int(\rho_n - \rho)^2 = \int\rho_n^2 - 2\int\rho_n\rho + \int\rho^2$, and each term converges (test $\rho_n^2$ against $1$ and $\rho_n$ against $\rho$). The argument with $\rho\log\rho$ is the same idea with a different convex function.
  \end{hint}
  \begin{solution}
    $\int(\rho_n - \rho)^2 = \int\rho_n^2 - 2\int\rho_n\rho + \int\rho^2$. Since $\Omega$ is bounded, $1 \in L^\infty$ and $\int\rho_n^2 \to \int\rho^2$. Since $\rho_n$ is bounded in $L^2$ and $\rho_n \weakto \rho$ in $L^1$, also $\rho_n \weakto \rho$ in $L^2$ (every subsequence has an $L^2$-weakly convergent subsequence, whose limit must be $\rho$), and $\rho \in L^2$; so $\int\rho_n\rho \to \int\rho^2$. Hence $\int(\rho_n - \rho)^2 \to \int\rho^2 - 2\int\rho^2 + \int\rho^2 = 0$.
  \end{solution}
\end{exercise}

\begin{exercise}
  The effective viscous flux in one dimension. For $(\rho u)_t + (\rho u^2)_x + p(\rho)_x = \nu u_{xx}$, show that $F = \nu u_x - p(\rho)$ satisfies $F_x = (\rho u)_t + (\rho u^2)_x$. Explain why this gives $F$ more regularity in $x$ than $u_x$ or $p(\rho)$ have separately.
  \begin{hint}
    Rearrange the momentum equation. The right-hand side is controlled by the energy bounds on $\rho u$ and $\rho u^2$, so $F$ gains a derivative, while $u_x$ and $p(\rho)$ alone can oscillate.
  \end{hint}
  \begin{solution}
    Rearranging the momentum equation, $\nu u_{xx} - p(\rho)_x = (\rho u)_t + (\rho u^2)_x$, and the left side is $(\nu u_x - p(\rho))_x = F_x$. The energy bounds control $\rho u$ and $\rho u^2$, so $F_x$ is a combination of a time derivative and a space derivative of bounded quantities, and $F$ itself is, roughly, an $x$-antiderivative of bounded quantities. By contrast, $\nu u_x$ is only in $L^2$ and $p(\rho)$ only in $L^1$, and each can oscillate in $x$; in $F$ their oscillations cancel. This extra regularity in $x$ is what lets $F$ pass to the limit in products such as $T(\rho_n)F_n$.
  \end{solution}
\end{exercise}

\subsection*{Where to find a proof}

\begin{itemize}
\item P.-L. Lions, \emph{Mathematical Topics in Fluid Mechanics, Vol. 2: Compressible Models}, Oxford, 1998.
\item E. Feireisl, A. Novotný, H. Petzeltová, \emph{J. Math. Fluid Mech.} 3 (2001) 358--392.
\item A. Novotný, I. Straškraba, \emph{Introduction to the Mathematical Theory of Compressible Flow}, Oxford, 2004, Chapter 7; E. Feireisl, \emph{Dynamics of Viscous Compressible Fluids}, Oxford, 2004.
\end{itemize}


\part{Case studies}
% One line per case study, in the order they should appear.
% To add one: copy sections/case-studies/_template.tex, fill it in, and add its \input line here.
% Chemotaxis compressible Navier–Stokes (vascular network formation)

\chapter{Chemotaxis compressible Navier--Stokes equations}\label{p:case-chemotaxis}
\status{draft}
\blurb{Huang, Rababah, Rawashdeh, J. Differential Equations 480 (2026). Local strong solutions with vacuum; global weak solutions for $\gamma > 3/2$.}

\begin{byline}Based on T. Huang, B. Rababah, A. Rawashdeh, \emph{On existence of chemotaxis compressible Navier--Stokes equations modeling vascular network formation}, J. Differential Equations 480 (2026) 114655. \href{https://doi.org/10.1016/j.jde.2026.114655}{doi:10.1016/j.jde.2026.114655}\end{byline}

This page walks through a complete existence theory for one nonlinear system. The aim is to show how the standard tools of PDE analysis fit together in a real argument. Every tool is linked to its own concept page, which states it precisely and points to a proof. Steps are tagged \routine{} when a standard tool is used in the standard way, and \newidea{} when the argument does something that cannot be taken off the shelf.

\conceptsused{%
  sobolev-spaces,
  bochner-spaces,
  weak-convergence,
  elementary-inequalities,
  sobolev-embedding,
  gagliardo-nirenberg,
  gronwall,
  elliptic-regularity,
  parabolic-regularity,
  transport-equation,
  notions-of-solution,
  energy-estimates,
  galerkin,
  uniqueness-energy,
  compactness,
  regularized-approximation,
  lions-feireisl}

\section{The model}

Endothelial cells, spread on a substrate, organize themselves into a network of blood vessels. They move like a compressible fluid and are attracted up the gradient of a chemical that they themselves secrete. In a bounded domain $\Omega \subset \R^3$ with smooth boundary, the unknowns are the cell density $\rho \ge 0$, the cell velocity $u \in \R^3$, and the chemoattractant concentration $\phi \ge 0$:

\begin{equation}\label{eq:system}
  \left\{
  \begin{aligned}
    &\rho_t + \dv(\rho u) = 0, \\
    &(\rho u)_t + \dv(\rho u \otimes u) + \nabla p(\rho) - \mu \Delta u - (\lambda+\mu)\nabla \dv u
       = \chi \rho \nabla \phi - \alpha \rho u, \\
    &\sigma \phi_t = D \Delta \phi + a \rho - b \phi,
  \end{aligned}
  \right.
  \end{equation}

with pressure $p(\rho) = \frac{\theta(\gamma-1)}{\gamma}\rho^\gamma$, $\gamma > 1$. The coefficients satisfy $D, \chi, \alpha, a, \theta > 0$, $\mu > 0$, $3\lambda + 2\mu > 0$, and $b, \sigma \ge 0$.

\begin{itemize}
\item $\chi \rho \nabla \phi$ is the chemotactic force: cells are pushed up the gradient of $\phi$.
\item $\alpha \rho u$ is friction with the substrate (damping).
\item $a\rho$ is secretion of the chemical by the cells; $b\phi$ is its decay.
\item The pressure models the fact that closely packed cells resist compression.
\end{itemize}

The boundary and initial conditions are

\begin{equation}\label{eq:bc}
    u = 0, \quad \frac{\partial \phi}{\partial \nu} = 0 \ \text{ on } \partial\Omega, \qquad
    (\rho, u, \phi)|_{t=0} = (\rho_0, u_0, \phi_0).
  \end{equation}

Setting $\mu = \lambda = 0$ gives the inviscid hyperbolic--parabolic model of Gamba et al. (2003), which reproduces the network patterns seen in experiments. Here we keep the viscosity.

\subsection*{Two difficulties to keep in mind}

\textbf{Vacuum.} The density may vanish on an open set. Where $\rho = 0$ the momentum equation degenerates: the term $\rho u_t$ gives no control of $u_t$. Both theorems below allow this.

\textbf{An energy with no sign.} The natural energy contains the cross term $-\chi\rho\phi$, which can be negative and large. Before any estimate can be used, one must show that the energy is bounded below.

\section{The main results}

\begin{thmbox}{Theorem A (local strong solutions)}
\index{chemotaxis Navier--Stokes!local strong solutions}%
Let $q \in (3,6]$. Suppose $\rho_0 \ge 0$, $\rho_0 \in W^{1,q}$, $u_0 \in D^2 \cap D^1_0$, $0 \le \phi_0 \in H^2$, and the compatibility condition

\begin{equation}\label{eq:compat}
      \mu \Delta u_0 + (\lambda+\mu)\nabla \dv u_0 - \nabla p(\rho_0) + \chi \rho_0 \nabla \phi_0 - \alpha \rho_0 u_0
      = \sqrt{\rho_0}\, g \quad \text{for some } g \in L^2(\Omega)
    \end{equation}

holds. Then there are $T_0 > 0$ and a unique strong solution of \eqref{eq:system}--\eqref{eq:bc} on $\Omega \times (0, T_0]$.
\end{thmbox}

\begin{proofslot}Proof of Theorem A: to be written. The map in Section 4 lists the steps in order.\end{proofslot}

\begin{thmbox}{Theorem B (global weak solutions)}
\index{chemotaxis Navier--Stokes!global weak solutions}%
Let $\gamma > 3/2$. Suppose $0 \le \rho_0 \in L^\gamma$ with $\rho_0 \not\equiv 0$, $\sqrt{\rho_0}\, u_0 \in L^2$, and $0 \le \phi_0 \in W^{2 - 2/\gamma,\, \gamma}$. Then for every $T > 0$ there is a finite-energy weak solution of \eqref{eq:system}--\eqref{eq:bc} on $\Omega \times (0,T)$.
\end{thmbox}

\begin{proofslot}Proof of Theorem B: to be written. The map in Section 5 lists the steps in order.\end{proofslot}

What "strong solution" and "finite-energy weak solution" mean, with the exact function spaces, is on \pref{notions-of-solution}{Strong, weak, and renormalized solutions}; the spaces themselves are on \pref{sobolev-spaces}{Sobolev spaces} and \pref{bochner-spaces}{spaces of time-dependent functions}.

\section{The energy structure}

Both theorems rest on one identity. Define

\begin{equation}\label{eq:energy}
    E(t) = \int_\Omega \Big( \tfrac12 \rho |u|^2 + \tfrac{\theta}{\gamma}\rho^\gamma
      + \tfrac{\chi D}{2a}|\nabla\phi|^2 + \tfrac{\chi b}{2a}\phi^2 - \chi\rho\phi \Big)\, dx .
  \end{equation}

\begin{step}{\routine{}The energy identity (Lemma 2.1)}
For smooth solutions, $\displaystyle \frac{d}{dt}E(t) = -\int_\Omega \Big( \tfrac{\chi\sigma}{a}|\phi_t|^2 + \alpha\rho|u|^2 + \mu|\nabla u|^2 + (\lambda+\mu)|\dv u|^2 \Big)dx \le 0.$

Test the momentum equation with $u$ and the $\phi$-equation with $\frac{\chi}{a}\phi_t$. The chemotactic work $\int \chi\rho\nabla\phi\cdot u$ is rewritten with the continuity equation as $\frac{d}{dt}\int\chi\rho\phi - \int\chi\rho\phi_t$, and the second piece cancels against the $\phi$-equation. Choosing the multipliers so that the cross terms cancel is the whole point.

\uses{\pref{energy-estimates}{energy estimates}.}
\end{step}

\begin{stepnew}{\newidea{}The energy is bounded below for $\gamma > 4/3$ (Proposition 2.2)}
Let $E_+$ be the energy without the cross term. Using only conservation of mass, $\int\rho = m$, one shows $\abs{\int\chi\rho\phi} \le C m^{q_0} + \frac12 E_+(t)$, so $E \ge \frac12 E_+ - Cm^{q_0}$. The cross term is split by Hölder, then $\norm{\rho}_{L^l}$ is interpolated between the conserved $L^1$ norm and the energy's $L^\gamma$ norm, and $\norm{\phi}_{L^{l'}}$ between $L^2$ and $H^1$. The exponents close (the power of $E_+$ stays below 1) exactly when $\gamma > 4/3$. Earlier work needed $\gamma > 8/5$. There are three cases, $\gamma > 2$, $\gamma = 2$, and $4/3 < \gamma < 2$; the last is where the exponent bookkeeping is delicate.

\uses{\pref{elementary-inequalities}{Hölder, Young, and L\textsuperscript{\textit{p}} interpolation}, \pref{gagliardo-nirenberg}{Gagliardo--Nirenberg}, \pref{sobolev-embedding}{Sobolev embeddings}.}
\end{stepnew}

\section{Map of the proof of Theorem A}

The strategy: first prove bounds that every smooth solution with $\rho_0 \ge \delta > 0$ must satisfy, with constants independent of $\delta$. Then build approximate solutions that satisfy the same bounds and pass to the limit. All the higher-order bounds are collected in one quantity,

\[
    \Phi(t) = \sup_{0 \le s \le t}\Big( \norm{\rho(s)}_{H^1 \cap W^{1,q}} + \norm{\nabla u(s)}_{L^2}
      + \norm{\sqrt{\rho}\,u_t(s)}_{L^2} + \norm{\phi(s)}_{H^2} + 1 \Big),
  \]

and the goal of Section 3 of the paper is the inequality $\Phi(t) \le \exp\big[C\mathcal{M} + C\int_0^t M(\Phi(s))\,ds\big]$, where $\mathcal{M}$ depends only on the data (through $\norm{g}_{L^2}$) and $M$ is some increasing function. Such an inequality forces $\Phi$ to stay bounded for a short time.

\begin{step}{\routine{}Step 1. Second derivatives of $u$ from the equation (Lemma 3.2)}
Read the momentum equation as $\Lame u = \rho u_t + \rho u\cdot\nabla u + \nabla p(\rho) - \chi\rho\nabla\phi + \alpha\rho u$ for the Lamé operator $\Lame$. Elliptic regularity bounds $\norm{\nabla u}_{H^1}$ by the right-hand side in $L^2$. The convection term produces $\norm{\nabla u}_{H^1}$ again, at a power below 2, so it is absorbed. Result: $\norm{\nabla u}_{H^1} \le M(\Phi)$.

\uses{\pref{elliptic-regularity}{elliptic regularity}, \pref{sobolev-embedding}{Sobolev embeddings}, \pref{gagliardo-nirenberg}{Gagliardo--Nirenberg}, \pref{elementary-inequalities}{Young's inequality}.}
\end{step}

\begin{step}{\routine{}Step 2. The time derivative $\sqrt\rho\, u_t$ (Lemmas 3.3--3.4)}
Differentiate the momentum equation in $t$, test with $u_t$, and estimate the six resulting terms. The value at $t = 0$, $\norm{\sqrt{\rho_0}\,u_t(0)}_{L^2}$, is exactly $\norm{g}_{L^2}$ from the compatibility condition \eqref{eq:compat}. This is why that condition is needed: without it the estimate has no starting value when $\rho_0$ vanishes somewhere.

\uses{\pref{energy-estimates}{energy estimates}, \pref{notions-of-solution}{compatibility conditions}, \pref{sobolev-embedding}{Sobolev embeddings}.}
\end{step}

\begin{step}{\routine{}Step 3. The density (Lemma 3.5)}
The continuity equation is a transport equation: it carries $\rho$ along the flow and never smooths it. Its $W^{1,q}$ norm grows at most like $\exp\big(C\int_0^t \norm{\nabla u}_{H^1 \cap D^{1,q}}\big)$. So one needs $\nabla u \in L^1(0,T; W^{1,q})$ with $q > 3$, which comes from the $W^{2,q}$ estimate for the Lamé system and Step 2.

\uses{\pref{transport-equation}{the transport equation}, \pref{elliptic-regularity}{elliptic regularity}.}
\end{step}

\begin{step}{\routine{}Step 4. The chemical (Lemma 3.6)}
Take the gradient of the $\phi$-equation, test with $\nabla\phi_t$, and use the Neumann condition. This bounds $\norm{\Delta\phi}_{L^2}$, hence $\norm{\phi}_{H^2}$ by elliptic regularity for the Neumann Laplacian.

\uses{\pref{parabolic-regularity}{parabolic regularity}, \pref{elliptic-regularity}{elliptic regularity}.}
\end{step}

\begin{step}{\routine{}Step 5. Closing the estimate}
Combine Steps 1--4 into the inequality for $\Phi$ above. A Gronwall-type continuity argument gives $T_0 > 0$, depending only on the data, with $\sup_{[0,T_0]}\Phi \le C\exp(C\norm{g}_{L^2})$.

\uses{\pref{gronwall}{Gronwall and continuity arguments}.}
\end{step}

\begin{step}{\routine{}Step 6. Construction (Section 4.1)}
Lift the density to $\rho_0 + \delta$ so there is no vacuum, adjusting $u_0$ so that the compatibility condition still holds with the same $g$. Project $u$ onto finitely many eigenfunctions of the Lamé operator; solve the continuity and $\phi$ equations exactly and the projected momentum equation as an ODE system (the mass matrix is invertible because $\rho \ge \delta > 0$). The bounds of Step 5 hold uniformly in the number of modes and in $\delta$, so one can pass to weak limits, first in the number of modes and then $\delta \to 0$.

\uses{\pref{galerkin}{the Faedo--Galerkin method}, \pref{weak-convergence}{weak and weak* convergence}, \pref{bochner-spaces}{spaces of time-dependent functions}.}
\end{step}

\begin{step}{\routine{}Step 7. Uniqueness (Section 4.2)}
Subtract two solutions. Estimate the difference in the lower norm $\norm{\sqrt{\rho_1}\,u}_{L^2}^2 + \norm{\rho}_{L^2}^2 + \norm{\phi}_{H^1}^2$, using the regularity of the two solutions only as coefficients. The result is a linear Gronwall inequality with zero initial value, so the difference vanishes.

\uses{\pref{uniqueness-energy}{uniqueness by difference estimates}, \pref{gronwall}{Gronwall's inequality}.}
\end{step}

\section{Map of the proof of Theorem B}

Weak solutions are built through a chain of regularized problems, each removing one difficulty, then the regularizations are removed one at a time. The paper concentrates on $3/2 < \gamma < 2$; larger $\gamma$ was covered by earlier work.

\begin{stepnew}{\newidea{}Step 1. The regularized system}
\[
    \left\{
    \begin{aligned}
      &\rho_t + \dv(\rho u) = \varepsilon\Delta\rho, \\
      &(\rho u)_t + \dv(\rho u\otimes u) + \nabla p(\rho) + \delta\nabla\rho^\beta + \varepsilon\nabla\rho\cdot\nabla u \\
      &\qquad = \mu\Delta u + (\lambda+\mu)\nabla\dv u + \chi\rho\nabla\phi - \alpha\rho u, \\
      &\sigma\phi_t + \varepsilon\Delta\big(|\Delta\phi|^{p_0}\Delta\phi\big) = D\Delta\phi + a\rho - b\phi,
      \qquad p_0 = \tfrac{2-\gamma}{\gamma-1}.
    \end{aligned}
    \right.
    \]

The first three regularizations are the standard Feireisl--Novotný--Petzeltová scheme: artificial viscosity $\varepsilon\Delta\rho$, a balancing term $\varepsilon\nabla\rho\cdot\nabla u$, and artificial pressure $\delta\rho^\beta$. The new difficulty: the artificial viscosity creates the term $-\varepsilon\int\phi\,\Delta\rho = -\varepsilon\int\rho\,\Delta\phi$ in the energy balance. Handled in the usual way, it leads to a term like $\varepsilon\int\rho^2$, which the energy controls only when $\gamma \ge 2$. The quasi-linear biharmonic term is designed for the range $3/2 < \gamma < 2$. It adds $\frac{\varepsilon\chi}{a(p_0+2)}\int|\Delta\phi|^{p_0+2}$ to the energy, and $p_0$ is chosen so that $p_0 + 2$ and $\gamma$ are conjugate exponents, so Young's inequality splits $\varepsilon\int\rho\,\Delta\phi$ between $\int\rho^\gamma$ and the new term.

\uses{\pref{regularized-approximation}{regularized approximation schemes}, \pref{elementary-inequalities}{Young's inequality}.}
\end{stepnew}

\begin{step}{\routine{}Step 2. Uniform energy bounds (Lemma 5.1)}
The energy identity of Section 3, redone for the regularized system, with the extra terms controlled as in Step 1 and the lower bound of Proposition 2.2. The constants do not depend on the Galerkin index, $\varepsilon$, or $\delta$.

\uses{\pref{energy-estimates}{energy estimates}, \pref{gronwall}{Gronwall's inequality}.}
\end{step}

\begin{step}{\routine{}Step 3. Removing the regularizations}
The Galerkin limit and $\varepsilon \to 0$ follow the compressible Navier--Stokes theory, using compactness to get strong convergence where the nonlinear terms need it. The biharmonic term disappears because the energy bounds $\varepsilon\int|\Delta\phi|^{p_0+2}$, so the term itself is $O(\varepsilon^{1/(p_0+2)})$ in the weak sense. The limit $\delta \to 0$ needs a little more integrability of $\rho$ than $L^\gamma$; this is where the pressure and the effective viscous flux enter.

\uses{\pref{compactness}{compactness}, \pref{lions-feireisl}{Lions--Feireisl theory}, \pref{weak-convergence}{weak convergence}.}
\end{step}

\begin{step}{\routine{}Step 4. Why $\gamma > 3/2$}
For $\rho\nabla\phi$ to make sense in the weak formulation, $\nabla\phi$ must lie in $L^{\gamma'}$, $\gamma' = \frac{\gamma}{\gamma-1}$, since $\rho$ is only in $L^\gamma$. Maximal regularity for the $\phi$-equation, with $a\rho \in L^\gamma$ as the source, gives $\phi \in W^{2,1}_\gamma$, so $\nabla\phi \in W^{1,\gamma}$ in space. The Sobolev embedding $W^{1,\gamma} \subset L^{3\gamma/(3-\gamma)}$ is enough when

\[
      \frac{3\gamma}{3-\gamma} > \frac{\gamma}{\gamma-1} \iff 3(\gamma - 1) > 3 - \gamma \iff \gamma > \tfrac32 .
    \]

The same threshold is where the Feireisl theory for compressible Navier--Stokes stops, so both obstructions meet at $3/2$.

\uses{\pref{parabolic-regularity}{maximal parabolic regularity}, \pref{sobolev-embedding}{Sobolev embeddings}.}
\end{step}

\section{Exercises}

These rebuild the paper's key computations by hand.

\begin{exercise}
  Derive the energy identity of Lemma 2.1 for smooth solutions: test the momentum equation with $u$ and the $\phi$-equation with $\frac{\chi}{a}\phi_t$, and show how the cross terms cancel. Check that the pressure potential is $\frac\theta\gamma\rho^\gamma$.
  \begin{hint}
    The chemotactic work: $\int\chi\rho\nabla\phi\cdot u = -\int\chi\phi\,\dv(\rho u) = \int\chi\phi\rho_t = \frac{d}{dt}\int\chi\rho\phi - \int\chi\rho\phi_t$. The $\phi$-equation tested with $\frac\chi a\phi_t$ produces $+\int\chi\rho\phi_t$. For the pressure, $P(\rho) = \frac{p(\rho)}{\gamma-1} = \frac\theta\gamma\rho^\gamma$ satisfies $\rho P'(\rho) - P(\rho) = p(\rho)$.
  \end{hint}
  \begin{solution}
    \emph{Momentum.} Test with $u$ ($u = 0$ on $\partial\Omega$). Using the continuity equation, $\int\big((\rho u)_t + \dv(\rho u\otimes u)\big)\cdot u = \frac{d}{dt}\int\frac12\rho|u|^2$. For the pressure, with $P(\rho) = \frac{\theta}{\gamma}\rho^\gamma$ one has $\rho P'(\rho) - P(\rho) = \theta\rho^\gamma - \frac\theta\gamma\rho^\gamma = p(\rho)$, and then $\int\nabla p\cdot u = -\int p\,\dv u = \frac{d}{dt}\int P(\rho)$. The viscous terms give $\mu\norm{\nabla u}^2 + (\lambda+\mu)\norm{\dv u}^2$. So

    \[ \frac{d}{dt}\int\Big(\frac12\rho|u|^2 + \frac\theta\gamma\rho^\gamma\Big) + \mu\norm{\nabla u}^2 + (\lambda+\mu)\norm{\dv u}^2
             = \int\chi\rho\nabla\phi\cdot u - \alpha\int\rho|u|^2. \]

    \emph{Chemotactic work.} $\int\chi\rho\nabla\phi\cdot u = -\int\chi\phi\,\dv(\rho u) = \int\chi\phi\rho_t = \frac{d}{dt}\int\chi\rho\phi - \int\chi\rho\phi_t$.

    \emph{Chemical.} Test $\sigma\phi_t = D\Delta\phi + a\rho - b\phi$ with $\frac\chi a\phi_t$, using $\partial_\nu\phi = 0$:

    \[ \frac{\chi\sigma}{a}\norm{\phi_t}^2 = -\frac{d}{dt}\int\Big(\frac{\chi D}{2a}|\nabla\phi|^2 + \frac{\chi b}{2a}\phi^2\Big) + \int\chi\rho\phi_t. \]

    Add the three: the term $\int\chi\rho\phi_t$ appears with opposite signs and cancels, leaving $\frac{d}{dt}E = -\int\big(\frac{\chi\sigma}{a}|\phi_t|^2 + \alpha\rho|u|^2 + \mu|\nabla u|^2 + (\lambda+\mu)|\dv u|^2\big)$.
  \end{solution}
\end{exercise}

\begin{exercise}
  Proposition 2.2 in its simplest case. Assume $b > 0$ and $\gamma > 2$. Using the mass $m$, show $\big|\int\chi\rho\phi\big| \le \frac12E_+ + Cm$, where $E_+$ contains $\frac\theta\gamma\int\rho^\gamma$ and $\frac{\chi b}{2a}\int\phi^2$.
  \begin{hint}
    $|\int\rho\phi| \le \norm{\rho}_{L^2}\norm{\phi}_{L^2} \le m^\eta\norm{\rho}_{L^\gamma}^{1-\eta}\norm{\phi}_{L^2}$ with $\eta = \frac{\gamma-2}{2(\gamma-1)}$ (Exercise 3 in \pref{elementary-inequalities}{Hölder, Young, and L\textsuperscript{\textit{p}} interpolation}). Use Young with three exponents: $\norm{\rho}_{L^\gamma}^{1-\eta}$ to the power $\frac{\gamma}{1-\eta}$, $\norm{\phi}_{L^2}$ to the power $2$, and $m^\eta$ to the power $r$ with $\frac{1-\eta}{\gamma} + \frac12 + \frac1r = 1$. You should find $r = \frac1\eta$, which needs $\gamma > 2$.
  \end{hint}
  \begin{solution}
    By Hölder and Exercise 3 in the inequalities section ($\rho \ge 0$, $\int\rho = m$): $\big|\int\rho\phi\big| \le \norm{\rho}_{L^2}\norm{\phi}_{L^2} \le m^\eta\norm{\rho}_{L^\gamma}^{1-\eta}\norm{\phi}_{L^2}$ with $\eta = \frac{\gamma-2}{2(\gamma-1)}$, so $1 - \eta = \frac{\gamma}{2(\gamma-1)}$.

    Young with three factors: $abc \le \frac{a^P}{P} + \frac{b^Q}{Q} + \frac{c^R}{R}$ when $\frac1P + \frac1Q + \frac1R = 1$. Take $a = \varepsilon_1\norm{\rho}_{L^\gamma}^{1-\eta}$ with $P = \frac{\gamma}{1-\eta} = 2(\gamma-1)$ (so $a^P = \varepsilon_1^P\norm{\rho}_{L^\gamma}^\gamma$), $b = \varepsilon_2\norm{\phi}_{L^2}$ with $Q = 2$, and $c = \frac{\chi m^\eta}{\varepsilon_1\varepsilon_2}$. Then $\frac1R = 1 - \frac{1}{2(\gamma-1)} - \frac12 = \frac{\gamma-2}{2(\gamma-1)} = \eta$, positive exactly when $\gamma > 2$, and $c^R = \big(\frac{\chi}{\varepsilon_1\varepsilon_2}\big)^{1/\eta}m$. Hence

    \[ \chi\Big|\int\rho\phi\Big| \le \frac{\varepsilon_1^P}{P}\int\rho^\gamma + \frac{\varepsilon_2^2}{2}\int\phi^2 + C(\varepsilon_1,\varepsilon_2)\,m. \]

    Choose $\varepsilon_1, \varepsilon_2$ with $\frac{\varepsilon_1^P}{P} \le \frac{\theta}{2\gamma}$ and $\frac{\varepsilon_2^2}{2} \le \frac{\chi b}{4a}$ (this uses $b > 0$). The first two terms are then at most half of $\frac\theta\gamma\int\rho^\gamma + \frac{\chi b}{2a}\int\phi^2 \le E_+$, so $\big|\int\chi\rho\phi\big| \le \frac12E_+ + Cm$.
  \end{solution}
\end{exercise}

\begin{exercise}
  The exponent $p_0$ in the regularized system. Show $p_0 + 2 = \frac{\gamma}{\gamma-1}$, the conjugate of $\gamma$. Then use Young to bound $\varepsilon\big|\int\rho\,\Delta\phi\big|$ by $\varepsilon\int\rho^\gamma$ and $\varepsilon\int|\Delta\phi|^{p_0+2}$.
  \begin{hint}
    $p_0 + 2 = \frac{2-\gamma+2\gamma-2}{\gamma-1}$. Then $\rho|\Delta\phi| \le \frac{\rho^\gamma}{\gamma} + \frac{|\Delta\phi|^{\gamma'}}{\gamma'}$.
  \end{hint}
  \begin{solution}
    $p_0 + 2 = \frac{2-\gamma}{\gamma-1} + 2 = \frac{2 - \gamma + 2\gamma - 2}{\gamma-1} = \frac{\gamma}{\gamma-1} = \gamma'$, and $\frac1\gamma + \frac1{\gamma'} = 1$. For $\rho \ge 0$, Young's inequality with exponents $\gamma, \gamma'$ (and a parameter $\kappa > 0$) gives

    \[ \varepsilon\Big|\int\rho\,\Delta\phi\Big| \le \varepsilon\int\rho|\Delta\phi| \le \varepsilon\kappa\int\rho^\gamma + \varepsilon C(\kappa)\int|\Delta\phi|^{p_0+2}. \]

    The regularized energy contains $\frac\theta\gamma\int\rho^\gamma$ and $\frac{\varepsilon\chi}{a(p_0+2)}\int|\Delta\phi|^{p_0+2}$, so both terms on the right are at most a constant times that energy. In the energy balance this gives $\frac{d}{dt}E \le CE$ (plus dissipation), and Gronwall closes the estimate. Without the biharmonic term one would need $\Delta\phi \in L^{\gamma'}$ from somewhere else, and $\gamma' > 2$ when $\gamma < 2$.
  \end{solution}
\end{exercise}

\section{Open questions}

\begin{itemize}
\item \textbf{Vanishing damping.} The construction of weak solutions never uses the damping term $\alpha\rho u$. What happens as $\alpha \to 0$?
\item \textbf{Vanishing viscosity.} Can weak solutions of the inviscid model ($\mu = \lambda = 0$) be obtained as limits of solutions of \eqref{eq:system}?
\item \textbf{The gap $4/3 < \gamma \le 3/2$.} The energy is bounded below there, but global weak solutions are not known to exist.
\item \textbf{Lower dimensions.} The same method should give weak solutions for $\gamma > 4/3$ in two dimensions and $\gamma > 1$ in one dimension.
\end{itemize}

\section*{References}

\begin{enumerate}[label={[\arabic*]}, leftmargin=*]
\item T. Huang, B. Rababah, A. Rawashdeh, On existence of chemotaxis compressible Navier--Stokes equations modeling vascular network formation, \emph{J. Differential Equations} 480 (2026) 114655.
\item X. Huo, A. Jüngel, Global existence and weak-strong uniqueness for chemotaxis compressible Navier--Stokes equations modeling vascular network formation, \emph{J. Math. Fluid Mech.} 26 (2024) 11.
\item Y. Cho, H. J. Choe, H. Kim, Unique solvability of the initial boundary value problems for compressible viscous fluids, \emph{J. Math. Pures Appl.} 83 (2004) 243--275.
\item E. Feireisl, A. Novotný, H. Petzeltová, On the existence of globally defined weak solutions to the Navier--Stokes equations, \emph{J. Math. Fluid Mech.} 3 (2001) 358--392.
\item A. Gamba et al., Percolation, morphogenesis, and Burgers dynamics in blood vessels formation, \emph{Phys. Rev. Lett.} 90 (2003) 118101.
\end{enumerate}


\appendix
\chapter{Hints to the exercises}
\printhints
\chapter{Solutions to the exercises}
\printsolutions

\end{document}
